\documentclass[11pt,a4paper]{article}
\usepackage[T1]{fontenc}
\usepackage{lmodern}
\usepackage[margin=25mm]{geometry}
\usepackage{amsmath,amssymb,amsthm,booktabs,longtable,array}
\usepackage[hidelinks]{hyperref}
\usepackage{xurl}
\usepackage{enumitem}
\setlist{itemsep=3pt,topsep=5pt}
\setlength{\emergencystretch}{2em}
\hypersetup{
    colorlinks=true,
    linkcolor=blue,
    urlcolor=blue,
    citecolor=blue,
    pdftitle={The Ninth-Dimensional Snake-in-the-Box Problem},
    pdfauthor={広中依緒璃, Tommy Ma, Steven Zhan}
}

\usepackage{microtype}
\newtheorem{theorem}{Theorem}[section]
\newtheorem{proposition}[theorem]{Proposition}
\newtheorem{lemma}[theorem]{Lemma}
\newtheorem{corollary}[theorem]{Corollary}
\newtheorem{conjecture}[theorem]{Conjecture}
\theoremstyle{definition}
\newtheorem{definition}[theorem]{Definition}
\theoremstyle{remark}
\newtheorem{remark}[theorem]{Remark}
\newcommand{\Q}{Q}
\newcommand{\Ham}{d_{\mathrm H}}
\title{Singleton-Square Bounds for the\\
Nine-Dimensional Snake-in-the-Box Problem}
\author{Iori Hironaka \and Tommy Ma \and Steven Zhan}
\date{\today}
\begin{document}
\maketitle
\begin{abstract}
Let $a(9)$ be the maximum number of edges in an induced path of the nine-dimensional binary hypercube. We establish the interval $191\le a(9)\le237$. The lower bound is the published construction of Orland et al., reproduced here as an explicit transition sequence. To our knowledge, the only explicit upper bound previously published for open induced paths in this dimension is the Danzer--Klee bound $2^{8}=256$. For the upper bound, we count coordinate squares meeting a path at exactly one vertex. A square-incidence identity converts lower bounds on these singleton counts into upper bounds on path length. We prove endpoint and adjacent-vertex restrictions, and exclude three consecutive internal vertices with singleton counts $2,2,2$. A finite potential argument then yields the bound 245. A survey of a local window model shows that window exclusions of this kind can raise the average singleton count to at least $3$. A small family of such window exclusions, each established by a formally verified replay of a refutation certificate, lowers the bound to 242 and then to 241. Direction-aware window automata, which also track the transition directions of the path, raise the average further; 175 window exclusions, established in the same formally verified way, give the bound 238, and 387 exclusions of a radius-four model give the bound 237. The construction and all upper bounds down to 237 have been checked in Lean under the same induced-path definition; the longest hand-written local classification is presented as a proof outline. A conditional bound 239 from a smaller automaton is subsumed by 238. Explicit adversarial configurations indicate that window arguments of this kind cannot reach much below about 227 edges. We also give an endpoint-maximality counterexample, exact conditions for joining paths in two layers, and explicit compatible paths of lengths 96 and 94. Restricted construction searches and a numerical conjecture are distinguished from the proved bounds. The exact value of $a(9)$ remains unresolved.
\end{abstract}
\clearpage
\setcounter{tocdepth}{1}
\tableofcontents
\clearpage

\section{Introduction and main results}
\begin{definition}
The graph $\Q_n$ has vertex set $\{0,1\}^n$, with adjacency $u\sim v$ when $\Ham(u,v)=1$. A snake of length $L$ is a sequence of distinct vertices $v_0,\ldots,v_L$ such that
\[
v_i\sim v_j\quad\Longleftrightarrow\quad |i-j|=1\qquad(i\ne j).
\]
Length counts edges. Thus a length-$L$ snake has $L+1$ vertices. Let $a(n)$ denote the largest possible length.
\end{definition}
Finiteness of $\Q_n$ ensures that a maximum exists. A transition sequence $t_1\cdots t_L$, with $0\le t_i<n$, specifies a path from the zero vertex by
\[
v_0=0,\qquad v_i=v_{i-1}\mathbin{\oplus}2^{t_i}.
\]
Integers encode binary words; $\oplus$ is exclusive-or. Translation by exclusive-or with a fixed word, coordinate permutations, and reversal preserve the snake conditions.

The exact values through dimension eight are listed in OEIS A099155 \cite{oeis}. Carlson and Hougen constructed a length-98 snake in dimension eight; the upper bound was established by exhaustive search by Östergård and Pettersson \cite{ostergard}. The universal bound $a(8)=98$ is used as a cited result here; a witness of that length is reproduced in Appendix~\ref{app:witness}, while the exhaustive upper bound is taken from that reference.

Orland et al. supply 191-edge snakes in dimension nine and a public collection of 131 inequivalent examples \cite{orland,data}. Optimality is unresolved. These examples are not known to exhaust all isometry classes of longest snakes. Its examples cannot be extended at their endpoints; that observation alone gives no maximum-length theorem.




\begin{theorem}[Main result]\label{thm:main}
The maximum length of an induced path in $\Q_9$ satisfies
\[
\boxed{191\le a(9)\le237.}
\]
\end{theorem}
The lower bound is verified from an explicit published sequence in Theorem~\ref{thm:lower}. The upper bound is proved in Theorem~\ref{thm:upper237}. It refines the bound 238 of Theorem~\ref{thm:upper238}, which refines the bound 241 of Theorem~\ref{thm:upper241}, which in turn refines the bound 245 of Theorem~\ref{thm:upper245}, whose main ingredients are the positivity of every singleton-square count, an endpoint count of at least four, the exclusion of adjacent internal counts $(1,1)$, $(1,2)$ and $(1,3)$, and the exclusion of three consecutive internal counts $(2,2,2)$. The bound 241 adds a small family of excluded count windows, the bound 238 a larger family of excluded windows that also records the transition directions, and the bound 237 such a family in a local model of larger radius (Section~\ref{sec:window-aggregate}). The earlier bounds 254, 251, 248, 247, 245, 242, 241 and 238 appear as intermediate consequences of the same counting identity.

The exposition separates three kinds of evidence. Explicit constructions and the upper-bound theorems are formally verified. The longest hand-written finite classification is described as a proof outline, with its precise counting definitions included. Solver outcomes are reported only for their stated restricted search spaces and are not premises of the upper-bound proof. The window exclusions behind the bounds 241, 238 and 237 were found with a solver, but the solver is not trusted: its refutations serve only as data for a proof checker whose soundness is itself formally proved. The conjectured value 196 arises from a separate numerical extrapolation and has no privileged role in the main theorem.

\subsection{Published bounds and related work}\label{sec:related}
In part of the literature cited below, \emph{snake} denotes an induced cycle. Here an induced cycle is called a \emph{coil}, \emph{snake} always means an open induced path, and lengths are counted in edges. For open snakes in $\Q_9$, the only explicit published upper bound that we have found is the bound $2^{n-1}$ of Danzer and Klee~\cite{danzerklee}, which gives 256 at $n=9$. Later upper bounds by Douglas~\cite{douglas}, Deimer~\cite{deimer}, Solov'eva~\cite{soloveva}, Emel'yanov~\cite{emelyanov}, Snevily~\cite{snevily}, Zémor~\cite{zemor}, Emelyanov and Lukito~\cite{emelyanovlukito}, and Lukito and van Zanten~\cite{lukito} concern coils, are asymptotic, or are aimed at larger dimensions; for instance, the abstract of~\cite{emelyanovlukito} states an improvement over earlier bounds for $37\le n\le19\,079$. As far as we could determine, none of them yields an explicit bound below 256 for open snakes in $\Q_9$. Bounds for coils do not transfer to open snakes without a separate argument: deleting a vertex of a coil gives a snake, but not conversely, and Theorem~\ref{thm:square-counterexample} below shows that a per-vertex square condition used for coils fails at internal vertices of open snakes. We have not been able to consult the full texts of several of these works, in particular \cite{soloveva,emelyanov,emelyanovlukito,lukito}, and we have not checked their exact statements; a small-dimension table there could contain a value for $n=9$. Subject to this caveat, we are not aware of an earlier published upper bound below 256 for open snakes in $\Q_9$, and to our knowledge the bounds below 256 proved here, culminating in 237, are new.

On the constructive side, Wynn obtained a 190-edge snake in $\Q_9$~\cite{wynn}, and Orland et al.\ improved this to 191 edges~\cite{orland}. The exact values $a(7)=50$ and $a(8)=98$ were obtained by exhaustive computer search; see~\cite{ostergard,oeis}.

\section{Exact targets for an optimality proof}
\begin{lemma}[Prefix restriction]\label{lem:prefix}
Every prefix of a snake is a snake.
\end{lemma}
\begin{proof}
Restrict the vertex sequence to its initial indices. Distinctness is inherited. Adjacency between any two retained indices has the same value and the same consecutive-index condition as in the original sequence.
\end{proof}

\begin{theorem}[Next-length exclusion]\label{thm:next}
For every integer $b\ge0$, all snakes in $\Q_9$ have length at most $b$ if and only if no snake has length $b+1$.
\end{theorem}
\begin{proof}
The upper bound directly excludes $b+1$. Conversely, a snake of any length larger than $b$ has a prefix of length $b+1$ by Lemma~\ref{lem:prefix}.
\end{proof}

\begin{corollary}
Conjecture~\ref{conj:196} is equivalent to the conjunction of a 196-edge construction and nonexistence of a 197-edge snake. A 196-edge construction would contain a 192-edge prefix and would refute the proposed upper bound 191.
\end{corollary}
\begin{proof}
Apply Theorem~\ref{thm:next} with $b=196$, and restrict a 196-edge snake to its first 193 vertices.
\end{proof}
For the competing optimum 191, the missing exclusion is the finite statement that no sequence $v_0,\ldots,v_{192}\in\{0,1\}^9$ is distinct and satisfies
\[
\Ham(v_i,v_j)=1\quad\Longleftrightarrow\quad |i-j|=1.
\]
The pairwise condition is required for every $0\le i<j\le192$. Neither next-length exclusion has been proved.

\section{Witnesses, endpoint obstruction, and an invalid inference}
\begin{theorem}[Verified lower bound]\label{thm:lower}
A 191-edge snake exists in $\Q_9$.
\end{theorem}
\begin{proof}
The first transition sequence in Appendix~\ref{app:witness} is the published Appendix E example of \cite{orland}. Starting at zero reconstructs 192 vertices. For each $0\le i<j\le191$, direct calculation gives $v_i\ne v_j$ and $\Ham(v_i,v_j)=1$ exactly when $j=i+1$. This is a finite verification of $\binom{192}{2}=18\,336$ unordered pairs using the displayed transition sequence, and establishes a snake of length 191.
\end{proof}

\begin{definition}
A snake $P=(v_0,\ldots,v_L)$ is blocked at its right endpoint when every neighbor $w$ of $v_L$ either equals a vertex $v_i$ with $i<L$, or is adjacent to such a vertex.
\end{definition}
\begin{lemma}[Blocking prevents endpoint extension]\label{lem:block}
A right-blocked snake cannot be extended by appending one vertex while retaining its entire ordered sequence.
\end{lemma}
\begin{proof}
An appended $w$ must be adjacent to $v_L$. If $w=v_i$ for an earlier index, distinctness fails. If $w\sim v_i$ for $i<L$, the appended vertex has a chord to an index at least two positions earlier. Both alternatives contradict the snake definition.
\end{proof}

\begin{theorem}[Endpoint maximality does not imply maximum length]
There is a 93-edge snake in $\Q_9$ blocked at both endpoints, although a 191-edge snake exists in the same cube.
\end{theorem}
\begin{proof}
The second sequence in Appendix~\ref{app:witness} specifies 94 vertices. Direct calculation verifies distinctness and the induced adjacency condition for all $\binom{94}{2}=4\,371$ unordered pairs. The appendix gives an earlier blocking vertex for each of the nine endpoint moves, in both the original and reversed orderings. Lemma~\ref{lem:block} excludes extension at either end. Theorem~\ref{thm:lower} supplies a strictly longer snake.
\end{proof}

\begin{theorem}[A sufficient universal blocking condition]
If every 191-edge snake in $\Q_9$ is blocked at its right endpoint, then $a(9)=191$.
\end{theorem}
\begin{proof}
If a 192-edge snake existed, its 191-edge prefix would admit extension by its remaining final vertex. This contradicts the universal blocking premise. Apply Theorem~\ref{thm:next} and the verified lower bound.
\end{proof}
The premise concerns every 191-edge snake. Blocking the 131 recorded classes is insufficient until completeness of the collection is established. Repeated recovery of the same classes is not a completeness proof.

\section{A global upper bound}
\begin{theorem}[Huang's degree theorem]\label{thm:huang}
Let $n\ge1$. If $H\subseteq\{0,1\}^n$ and $|H|>2^{n-1}$, then the subgraph induced by $H$ has a vertex of degree at least $\sqrt n$.
\end{theorem}
\begin{proof}[Proof argument]
Define real symmetric matrices recursively by $A_0=(0)$ and
\[
A_n=\begin{pmatrix}A_{n-1}&I\\I&-A_{n-1}\end{pmatrix}.
\]
Induction gives $A_n^2=nI$. Its nonzero entries are precisely signed hypercube edges. For $n\ge1$, its trace is zero, so its eigenvalues $+\sqrt n$ and $-\sqrt n$ each have multiplicity $2^{n-1}$. The subspace of vectors supported on $H$ has dimension $|H|$. If $|H|>2^{n-1}$, it has nonzero intersection with the positive eigenspace, by the dimension inequality for intersections. Choose a nonzero vector $y$ in that intersection and $q\in H$ maximizing $|y_q|>0$. Then
\[
\sqrt n\,|y_q|=|(A_ny)_q|
\le\sum_{p\in H,\ p\sim q}|y_p|
\le \deg_H(q)|y_q|.
\]
Division proves the assertion. This is Huang's argument \cite{huang}, included to make the dependence of the upper bound explicit.
\end{proof}

\begin{lemma}[Compatibility of the integer and bit-function encodings]
Map $v\in\{0,\ldots,511\}$ to its nine binary digits $E(v)(k)$. This map is injective. If $E(u)$ and $E(v)$ differ in exactly one coordinate, then $v=u\oplus2^k$ for some $k<9$.
\end{lemma}
\begin{proof}
All digits at indices at least nine vanish because $v<2^9$. Equality of the first nine digits therefore gives equality at every digit and hence equality of integers. If exactly digit $k$ differs, exclusive-or with $2^k$ changes that digit and preserves every other digit. Equality of all digits then gives the stated integer identity.
\end{proof}

\begin{lemma}[Degree bound for an encoded snake]
The set of bit-function encodings of a snake's vertices induces maximum degree at most two.
\end{lemma}
\begin{proof}
For a vertex with index $i$, any encoded neighbor corresponds, by the preceding lemma, to an integer-encoded adjacent vertex. The snake condition forces its index to be $i-1$ or $i+1$, when those indices exist. There are at most two such vertices. Distinctness and injectivity of $E$ preserve the vertex count.
\end{proof}

\begin{theorem}[Verified universal bound]\label{thm:upper}
Every snake in $\Q_9$ has at most 255 edges. Consequently,
\[
\boxed{191\le a(9)\le255.}
\]
\end{theorem}
\begin{proof}
A snake of at least 256 edges would have at least 257 vertices, exceeding $2^8$. Theorem~\ref{thm:huang} would give an induced degree of at least $\sqrt9=3$, contradicting the preceding degree bound. Combine the upper bound with Theorem~\ref{thm:lower}.
\end{proof}
The bound is a standard consequence of Huang's theorem. The square-incidence arguments below give a stronger bound.


\section{Square faces and the role of endpoints}
For a vertex set $S\subseteq\Q_9$ and $x\in S$, let $\alpha_S(x)$ be the number of two-dimensional coordinate faces whose intersection with $S$ is exactly $\{x\}$. Such a face will be called a singleton square at $x$. Square-face counting occurs in the cycle bounds reviewed by Lukito and van Zanten \cite{lukito}. Their cycle argument uses a lower bound of two singleton squares at each selected vertex. A direct extension of that condition to every internal vertex of every open snake is false.

\begin{theorem}[An internal-vertex counterexample]\label{thm:square-counterexample}
There is a 137-edge snake $P=(v_0,\ldots,v_{137})$ in $\Q_9$ with $v_{127}=0$ and $\alpha_{V(P)}(0)=1$. The vertex $v_{127}$ is at least ten path edges from either endpoint.
\end{theorem}
\begin{proof}
Use the transition sequence in Appendix~\ref{app:square-counterexample}, starting at the integer vertex 145. Direct calculation checks all $\binom{138}{2}=9\,453$ unordered vertex pairs: the vertices are distinct, and adjacency holds precisely for consecutive indices. It also gives
\[
v_0=145,\qquad v_{127}=0,\qquad v_{137}=288.
\]
The selected neighbors of zero are 1 and 2. Every square at zero involving coordinate 0 or 1 therefore contains another selected vertex. For $2\le a<b\le8$, neither $2^a$ nor $2^b$ is selected, and the opposite vertex $2^a+2^b$ is selected except when $(a,b)=(2,3)$. Thus the unique singleton square is $\{0,4,8,12\}$. Finally, the path distances from $v_{127}$ to the endpoints are 127 and 10.
\end{proof}
The Hamming distances from zero to the two endpoints are 3 and 2. Distance along the path is therefore insufficient for ignoring endpoint effects in this local cube configuration. This example concerns a 137-edge path; it makes no claim about the corresponding local condition for all longest paths.

\begin{lemma}[At least one singleton square]\label{lem:one-square}
If $S$ is the vertex set of an induced path in $\Q_9$, then $\alpha_S(x)\ge1$ for every $x\in S$.
\end{lemma}
\begin{proof}
Translate $x$ to zero. Translation preserves adjacency and the induced-path property. Addition below is coordinatewise modulo two, and $e_i$ denotes the unit vector in coordinate $i$.

First suppose that the path has at most two edges. The three coordinate squares using the pairs $(0,1)$, $(2,3)$, and $(4,5)$ have mutually disjoint sets of nonzero vertices. If none were a singleton square, they would contribute three distinct selected vertices besides zero, exceeding the number of vertices in the path. Thus at least one is a singleton square.

Now suppose that the path has at least three edges, and assume for contradiction that no singleton square exists at zero. For any two distinct coordinates $i,j$ whose unit vectors are absent, the opposite vertex $e_i+e_j$ must then be selected.

Every vertex of a path with at least three edges has an internal path vertex as a neighbor. Choose such a neighbor $e_a$ of zero, and let $c$ be its other path neighbor. Write $c=e_a+e_i$, where $i\ne a$ because $c\ne0$. The unit vector $e_i$ is absent: otherwise $0,e_a,c,e_i,0$ would be a selected four-cycle, which an induced path cannot contain.

At most two unit vectors are selected, since zero has induced degree at most two. Choose a coordinate $b\notin\{a,i\}$ so that every selected unit vector is indexed by $a$ or $b$. If there is a second selected unit vector, take its coordinate; otherwise any coordinate outside $\{a,i\}$ suffices. Set
\[
U=\{0,\ldots,8\}\setminus\{a,b\}.
\]
Then $i\in U$, every unit vector indexed by $U$ is absent, and every pair vertex $e_j+e_k$ with distinct $j,k\in U$ is selected. Every triple vertex $e_j+e_k+e_l$ with distinct $j,k,l\in U$ is absent, because its three pair-vertex neighbors are selected and a path has maximum degree two.

Choose four distinct coordinates $j\in U\setminus\{i\}$; six are available. For each of these four coordinates, define
\[
p_j=e_i+e_j,\qquad q_j=e_a+e_i+e_j.
\]
All four $p_j$ are selected. The four $q_j$ are distinct neighbors of $c$, distinct also from its already selected neighbor $e_a$. Therefore at most one $q_j$ is selected, and at least three are absent.

For each absent $q_j$, the corresponding $p_j$ has at most one selected neighbor. Flipping coordinate $i$ or $j$ gives an absent unit vector; flipping any other coordinate in $U$ gives an excluded triple; and flipping $a$ gives the absent $q_j$. Only the neighbor $e_b+e_i+e_j$ can remain selected. Hence at least three distinct selected vertices have induced degree at most one. In a path with at least three edges, precisely the two endpoints have degree at most one. This contradiction proves the lemma.
\end{proof}

\begin{theorem}[Square-face refinement]\label{thm:square-upper}
Every snake in $\Q_9$ has at most 254 edges. Together with the explicit witness, this gives
\[
\boxed{191\le a(9)\le254.}
\]
\end{theorem}
\begin{proof}
A path with one vertex satisfies the assertion, so let $S$ be a path vertex set of size $v\ge2$. For $0\le a<b\le8$ and $x\in\Q_9$, write
\[
C_{ab}(x)=\{x,x+e_a,x+e_b,x+e_a+e_b\},
\]
where addition is coordinatewise modulo two. Define three indicators, each taking values in $\{0,1\}$:
\begin{align*}
s_{ab}(x)&=1\quad\Longleftrightarrow\quad C_{ab}(x)\cap S=\{x\},\\
w_{ab}(x)&=1\quad\Longleftrightarrow\quad x,x+e_a,x+e_b\in S,\\
e_{ab}(x)&=1\quad\Longleftrightarrow\quad C_{ab}(x)\cap S=\varnothing.
\end{align*}
All three indicators are zero when their stated condition fails. Let
\[
T=\sum_{a<b}\sum_{x\in\Q_9}s_{ab}(x),\qquad
W=\sum_{a<b}\sum_{x\in\Q_9}w_{ab}(x),\qquad
E=\sum_{a<b}\sum_{x\in\Q_9}e_{ab}(x).
\]
No square has all four vertices selected, since that would be a four-cycle in the induced path. For each square, considering its possible numbers of selected vertices, zero through three, gives
\begin{equation}\label{eq:local-square-identity}
2e_{ab}(x)+\sum_{y\in C_{ab}(x)}s_{ab}(y)
 +\sum_{y\in C_{ab}(x)}\mathbf{1}_S(y)
=2+\sum_{y\in C_{ab}(x)}w_{ab}(y).
\end{equation}
For three selected vertices, exactly one is adjacent to the other two, so the last sum is one. For zero, one, or two selected vertices, the last sum is zero; the other terms give the asserted identity directly.

For each fixed coordinate pair, the four translations taking $x$ to the corners of $C_{ab}(x)$ are permutations of $\Q_9$. Summing~\eqref{eq:local-square-identity} over all 512 choices of $x$ and all 36 coordinate pairs therefore gives
\begin{equation}\label{eq:anchored-square-count}
2E+4T+144v=36\,864+4W.
\end{equation}
Here every selected vertex contributes four times for each coordinate pair to the vertex-incidence sum. The quantity $E$ counts each empty square once from each of its four corners, so it is four times the number of distinct empty squares.

Lemma~\ref{lem:one-square} implies $T\ge v$. A nonzero contribution to $W$ requires two distinct selected neighbors of its selected center, so the center is an internal path vertex. Since an induced path has degree at most two, at most one unordered coordinate pair contributes at any such center. There are $v-2$ internal vertices, hence $W\le v-2$. Using $E\ge0$ in~\eqref{eq:anchored-square-count},
\[
4v+144v\le36\,864+4(v-2),
\qquad\text{so}\qquad144v\le36\,856<144\cdot256.
\]
Thus $v\le255$, and the number of edges is $v-1\le254$.
\end{proof}


\section{Adjacent singleton minima and the bound 251}
This section strengthens the preceding incidence estimate by restricting how vertices with exactly one singleton square can occur along a path. The 137-edge counterexample still shows that such a vertex can occur internally. The additional restriction concerns two adjacent internal vertices, rather than either vertex in isolation.

Fix an induced path with vertex set $S$. In this section, call a vertex selected when it belongs to $S$, and abbreviate $e_a$ by $a$, $e_a+e_b$ by $ab$, and similarly for longer coordinate sums. All additions are in $\mathbb F_2^9$. Within a missing-pair family, the abbreviation $ij$ instead denotes the unordered coordinate pair $\{i,j\}$; in a vertex expression it denotes the sum of the corresponding unit vectors. Coordinate letters in a displayed sum are distinct unless stated otherwise. Every selected vertex has at most two selected neighbors, and the only selected vertices with at most one selected neighbor are the two endpoints, provided the path has an edge. The selected graph contains no cycle.

\begin{lemma}[Restriction on a continuation]\label{lem:singleton-arm}
Suppose that zero is an internal selected vertex with selected neighbors $a,b$, and that its unique singleton square uses coordinates $c,d$. Then $a,b,c,d$ are distinct. If $ai$ is selected and $i\ne a$, then $i\in\{c,d\}$. The analogous assertion holds with $a,b$ interchanged.
\end{lemma}
\begin{proof}
A square involving coordinate $a$ or $b$ is not a singleton at zero, so $c,d$ avoid $a,b$. Put $U=\{0,\ldots,8\}\setminus\{a,b\}$. The unit vertices indexed by $U$ are absent, and every pair vertex indexed by two distinct elements of $U$ is selected except $cd$.

If $i=b$, the vertices $0,a,ab,b$ form a selected four-cycle. Thus a selected continuation $ai\ne0$ has $i\in U$. Suppose that $i\notin\{c,d\}$. There are four coordinates $j$ outside $\{a,b,i,c,d\}$; choose all four. Each $ij$ is selected. For any $k\in U\setminus\{i,j\}$, all three pair vertices $ij,ik,jk$ are selected: neither $i$ nor $j$ is $c$ or $d$, so none of these pairs can be the exceptional pair $cd$. Hence $ijk$ is absent, since otherwise it would have three selected neighbors.

Consequently, the only possible selected neighbors of $ij$ are $aij$ and $bij$. The four vertices $aij$ are distinct neighbors of the selected vertex $ai$, which already has the selected neighbor $a$. At most one $aij$ is selected. For at least three of the chosen coordinates $j$, the selected vertex $ij$ therefore has at most one selected neighbor. These are three distinct endpoints, a contradiction. This also covers a selected neighbor $a$ that might initially have been an endpoint: the existence of a further selected neighbor $ai\ne0$ itself makes $a$ internal.
\end{proof}

\begin{theorem}[No adjacent internal singleton minima]\label{thm:no-adjacent-singletons}
If $x,y$ are adjacent internal vertices of an induced path in $\Q_9$, then they cannot both satisfy $\alpha_S(x)=\alpha_S(y)=1$.
\end{theorem}
\begin{proof}
Assume otherwise and translate $x$ to zero. Write $y=a$, let $b$ be the other selected neighbor of zero, and let $ac$ be the other selected neighbor of $a$. The coordinates $a,b,c$ are distinct: equality $a=c$ repeats zero, while $b=c$ creates the selected four-cycle $0,a,ab,b$.

By Lemma~\ref{lem:singleton-arm}, the unique singleton square at zero uses a pair $\{c,d\}$, where $d$ avoids $a,b,c$. Apply the same lemma after translating $a$ to zero. Its selected neighbors then have coordinates $a,c$, and the continuation through the original zero to $b$ has second coordinate $b$. Thus its unique singleton square uses $\{b,e\}$, where $e$ avoids $a,b,c$. The coordinates $d,e$ may coincide. The two singleton-square hypotheses give the following complete pair information:
\begin{itemize}
\item Every pair $ij$ with $i,j\notin\{a,b\}$ is selected except $cd$.
\item Every vertex $aij$ with $i,j\notin\{a,c\}$ is selected except $abe$.
\end{itemize}

Let $R$ be the coordinates outside $\{a,b,c,d,e\}$. There are at least four such coordinates. For $k\in R$, set $t_k=abk$, which is selected by the second pair statement. We first classify all possible selected neighbors of $t_k$.

The vertices $ak$ and $ab$ are absent, since the selected vertex $a$ already has selected neighbors $0,ac$. The vertex $bk$ is absent by Lemma~\ref{lem:singleton-arm} at zero. The same lemma at the translated vertex $a$ shows that $ack$ is absent. These conclusions do not assume that $b$ or $ac$ is internal: a selected continuation would make the corresponding vertex internal, so the lemma still rules it out.

For any coordinate $l\notin\{a,b,c,e,k\}$, the three vertices $abk,abl,akl$ are selected. They are distinct neighbors of $abkl$, which must therefore be absent. Enumerating the nine possible coordinate flips now shows that the only possible selected neighbors of $t_k$ are
\[
abck\quad\hbox{and}\quad abek.
\]
In particular, if $t_k$ is internal, both are selected. Whenever $abek$ is selected, the vertex $bek$ is absent: otherwise $abek$ would have the three distinct selected neighbors $abk,aek,bek$.

\paragraph{Case $e\ne d$.}
For each $k\in R$, the pair $ek$ is selected. Choose three distinct coordinates $k$ in $R$. The three sets $\{t_k,ek\}$ are pairwise disjoint: every vertex in one set contains its corresponding coordinate $k$ and contains neither of the other two chosen coordinates. Since a path has only two endpoints, at least one of these sets consists entirely of internal vertices. Fix its coordinate $k$.

The internal vertex $t_k$ forces $abek$ to be selected and hence $bek$ to be absent. We claim that the selected vertex $ek$ can then have only the selected neighbor $aek$. Flipping $e$ or $k$ gives an absent unit vertex. Flipping $b$ gives the absent $bek$. For any other coordinate $l\notin\{a,b,e,k\}$, the three pair vertices $ek,el,kl$ are selected: neither $e$ nor $k$ belongs to the exceptional pair $\{c,d\}$. Thus $ekl$ is absent because it would have three selected neighbors. Only the flip in coordinate $a$ remains. This makes $ek$ an endpoint, contradicting its choice as internal.

\paragraph{Case $e=d$.}
For $k\in R$, the vertices $t_k,ck,dk$ are all selected. Choose three distinct coordinates in $R$. The corresponding three sets $\{t_k,ck,dk\}$ are pairwise disjoint by the same coordinate argument. One of these sets therefore contains no endpoint. Fix its coordinate $k$.

Since $t_k$ is internal, both $abck$ and $abdk$ are selected, and the preceding degree argument makes $bdk$ absent. The only possible selected neighbors of $dk$ are now $adk$ and $cdk$. Indeed, its unit neighbors are absent, $bdk$ is absent, and any triple $dkl$ with $l\notin\{a,b,c,d,k\}$ has the three selected pair neighbors $dk,dl,kl$. Since $dk$ is internal, $cdk$ is selected.

Likewise, the only possible selected neighbors of $ck$ are $bck$ and $cdk$: the candidate $ack$ is absent by the continuation lemma, and every remaining triple outside the exceptional coordinate $d$ has three selected pair neighbors. Since $ck$ is internal, $bck$ is selected. The second pair statement also gives the selected vertex $adk$. We have therefore obtained the selected eight-cycle
\[
abk,\ abck,\ bck,\ ck,\ cdk,\ dk,\ adk,\ abdk,\ abk.
\]
Its eight vertices are distinct because $a,b,c,d,k$ are distinct, and every displayed edge flips one coordinate. A path cannot contain this cycle. Both cases are contradictory, proving the theorem.
\end{proof}

\begin{lemma}[Aggregate singleton count]\label{lem:aggregate-singleton}
For a snake with $m\ge1$ edges, let $T=\sum_{i=0}^{m}\alpha_S(v_i)$. Then $2T\ge3m$.
\end{lemma}
\begin{proof}
Write $\alpha_i=\alpha_S(v_i)$. Lemma~\ref{lem:one-square} gives $\alpha_i\ge1$. If $m=1$, then $2T\ge4\ge3m$. For $m\ge2$, each pair of adjacent internal vertices has singleton counts summing to at least three by Theorem~\ref{thm:no-adjacent-singletons}. The first and last adjacent pairs each have sum at least two. Summing all $m$ adjacent pairs gives
\[
2T-\alpha_0-\alpha_m
=\sum_{i=0}^{m-1}(\alpha_i+\alpha_{i+1})
\ge 3m-2.
\]
Adding $\alpha_0+\alpha_m\ge2$ proves $2T\ge3m$.
\end{proof}

\begin{theorem}[Verified aggregate refinement]\label{thm:upper251}
Every snake in $\Q_9$ has at most 251 edges. In particular,
\[
\boxed{191\le a(9)\le251.}
\]
\end{theorem}
\begin{proof}
A one-vertex path is immediate. Otherwise let $v=m+1\ge2$ be the vertex count. The anchored incidence identity~\eqref{eq:anchored-square-count}, together with $E\ge0$ and $W\le v-2$, gives
\[
140v+4T\le36\,856.
\]
Lemma~\ref{lem:aggregate-singleton} gives $4T\ge6(v-1)$, so
\[
146v\le36\,862<146\cdot253.
\]
Thus $v\le252$, hence $m\le251$. Combine this with Theorem~\ref{thm:lower}.
\end{proof}



\section{Adjacent singleton counts one and two: the bound 248}
We retain the selected-vertex and coordinate-sum notation of the preceding section. A selected vertex with at most one selected neighbor will be called an endpoint in the argument below. There are at most two such vertices in an induced path, including the degenerate one-vertex case.

\begin{theorem}[Exclusion of adjacent counts one and two]\label{thm:no-adjacent-one-two}
If $x,y$ are adjacent internal vertices of an induced path in $\Q_9$, then $\alpha_S(x)=1$ and $\alpha_S(y)=2$ cannot hold simultaneously.
\end{theorem}
\begin{proof}
Translate $x$ to zero and write $y=a$. As before, let $b$ be the other selected neighbor of zero and $ac$ the other selected neighbor of $a$. Lemma~\ref{lem:singleton-arm} shows that the unique missing pair at zero is $cd$, with $a,b,c,d$ distinct. Put
\[
U=\{0,\ldots,8\}\setminus\{a,b,c,d\},\qquad |U|=5.
\]
Let $F$ be the two missing unordered coordinate pairs at $a$. Since the selected neighbors of $a$ have directions $a,c$, each member of $F$ is a two-element subset of $\{b,d\}\cup U$. For distinct $i,j$ outside $\{a,c\}$, the vertex $aij$ is selected if and only if $\{i,j\}\notin F$. This is exactly the singleton-square definition after translation by $a$; the two missing pairs are distinct.

Write $t_i=abi$ for $i\in U$, and also $t_d=abd$. Every pair vertex $ij$ with $i,j$ outside $\{a,b\}$ is selected except $cd$. For $i\in U$, the vertices $ai,ab$ are absent because $a$ already has selected neighbors $0,ac$, while $bi$ is absent by Lemma~\ref{lem:singleton-arm}. Thus the only possible selected neighbors of $t_i$ are $abci$ and
\[
abij\qquad\bigl(j\in\{d\}\cup(U\setminus\{i\})\bigr).
\]
For distinct $i,j\in\{d\}\cup U$ with at least one index in $U$, the selected pair $ij$ can have selected neighbors only $aij,bij$ and, when one index is $d$, possibly $cij$. To see this, its unit neighbors are absent. Any other triple neighbor has three selected pair neighbors and is therefore absent; the only possible failure of this three-neighbor argument involves the exceptional missing pair $cd$.

We use four consequences of these observations.
\begin{enumerate}
\item If $t_i,t_j,aij$ are selected, then $abij$ is absent, because it would have these three distinct selected neighbors.
\item If $t_i,t_j,abij$ are selected and $\{i,j\}\in F$, then $ij$ is an endpoint. Indeed, $aij$ is absent by membership in $F$, and $bij$ is absent because $abij$ already has selected neighbors $t_i,t_j$. The pair-neighbor classification leaves at most the single candidate $cij$.
\item Suppose $h,i$ are distinct elements of $U$, $\{b,h\}\in F$, $\{h,i\}\notin F$, and $t_i,abhi$ are selected. Then $hi$ is an endpoint. The vertex $ahi$ is selected, so the two selected neighbors $t_i,ahi$ of $abhi$ force $bhi$ to be absent. The only possible selected neighbor of $hi$ is consequently $ahi$.
\item Suppose $i\in U$, both $t_i,adi$ are selected, and the only possible selected neighbors of $t_i$ are $abci,abdi$. Then at least one of $t_i,ci,di$ is an endpoint. Otherwise both quads are selected. The vertex $aci$ is absent, since it would have the three selected neighbors $ac,ci,abci$. Also $bdi$ is absent, since $abdi$ already has selected neighbors $t_i,adi$. The internal pair $di$ then forces $cdi$ to be selected, and the internal pair $ci$ forces $bci$ to be selected. This gives the impossible selected eight-cycle
\[
t_i,\ abci,\ bci,\ ci,\ cdi,\ di,\ adi,\ abdi,\ t_i.
\]
This last observation uses no singleton-count-one assumption at $a$.
\end{enumerate}

\paragraph{Case 1: neither member of $F$ contains $b$.}
All five $t_i$ and $t_d$ are selected. A non-endpoint $t_i$ needs a selected neighbor other than $abci$, hence some $abij$. Consequence 1 forces $\{i,j\}\in F$, and consequence 2 makes $ij$ an endpoint.

Let $E$ consist of those pairs in $F$ whose pair vertices are endpoints, and set $k=|E|\le2$. These $k$ distinct pair vertices differ from all $t_i$, whose supports have size three. At most $2-k$ of the five $t_i$ are endpoints, so at least $3+k$ are not. Each of these indices lies in a member of $E$, but $k$ pairs cover at most $2k$ indices. Thus $3+k\le2k$, impossible for $k\le2$.

\paragraph{Case 2: $\{b,d\}\in F$.}
First suppose the other missing pair is $\{f,g\}$ and does not contain $b$. At least three indices $i\in U$ avoid $f,g$. For each such index, $t_i,adi$ are selected. Every possible $abij$ with $j\in U\setminus\{i\}$ is absent by consequence 1, since $t_i,t_j,aij$ are selected. Only $abci,abdi$ remain as possible selected neighbors. Consequence 4 gives an endpoint in each of the three sets $\{t_i,ci,di\}$.

These sets are pairwise disjoint. Indeed, every vertex in the set indexed by $i$ contains coordinate $i$ and contains none of the other chosen indices. This would give at least three endpoints.

Otherwise $F=\{\{b,d\},\{b,h\}\}$ with $h\in U$. For each of the four $i\in U\setminus\{h\}$, the only possible selected neighbors of $t_i$ are $abci,abdi,abhi$, and $t_i,adi$ are selected. If $abhi$ is selected, consequence 3 makes $hi$ an endpoint. If it is absent, consequence 4 supplies an endpoint among $t_i,ci,di$. The four sets $\{t_i,ci,di,hi\}$ are pairwise disjoint by the same coordinate argument, again impossible.

\paragraph{Case 3: $\{b,d\}\notin F$, but $\{b,e\}\in F$ for some $e\in U$.}
If the other missing pair is $\{b,h\}$, with $h\in U\setminus\{e\}$, consider the three indices $i\in U\setminus\{e,h\}$. The vertices $t_d,t_i,adi$ are selected, so consequence 1 excludes $abdi$. Every possible selected neighbor of $t_i$ other than $abci,abei,abhi$ is similarly excluded. If $t_i$ is not an endpoint, at least one of $abei,abhi$ is selected, and consequence 3 makes respectively $ei$ or $hi$ an endpoint. The three sets $\{t_i,ei,hi\}$ are pairwise disjoint, a contradiction.

It remains that the other missing pair is $\{f,g\}$ and does not contain $b$. Put $U_0=U\setminus\{e\}$. Any $i\in U_0\setminus\{f,g\}$ has $t_i,t_d,adi,aei$ selected. Consequence 1 excludes $abdi$ and every other possible selected neighbor of $t_i$ except $abci,abei$. Hence either $t_i$ is an endpoint or $abei$ is selected; in the latter event consequence 3 makes $ei$ an endpoint. If at most one of $f,g$ lies in $U_0$, there are at least three pairwise disjoint endpoint-containing sets $\{t_i,ei\}$.

The final possibility has distinct $f,g\in U_0$. The other two indices of $U_0$ already give two disjoint endpoint-containing sets. If $abfg$ is selected, consequence 2 gives the third endpoint $fg$, distinct from both sets. If $abfg$ is absent, the same two-neighbor argument applies also to $f$ and $g$: the only additional potential quad for either index was $abfg$. Consequently all four pairwise disjoint sets $\{t_i,ei\}$, for $i\in U_0$, contain endpoints. Both alternatives are impossible.

The three cases exhaust every two-element missing-pair set $F$, proving the theorem. Reversing the orientation of the adjacent pair also excludes counts $(2,1)$.
\end{proof}

\begin{lemma}[Strengthened aggregate count]\label{lem:aggregate-four}
For a snake with $m\ge1$ edges and total singleton count $T$, one has $4m\le2T+2$.
\end{lemma}
\begin{proof}
Every singleton count is a positive integer. Theorems~\ref{thm:no-adjacent-singletons} and~\ref{thm:no-adjacent-one-two} therefore imply that the counts of any two adjacent internal vertices sum to at least four. For $m\ge2$, summing the $m$ adjacent-pair counts, with the first and last pairs bounded below by two, gives
\[
2T-\alpha_0-\alpha_m\ge4m-4.
\]
Since $\alpha_0+\alpha_m\ge2$, this proves $2T\ge4m-2$. If $m=1$, the universal lower bound $T\ge2$ proves the assertion directly.
\end{proof}

\begin{theorem}[Verified upper bound 248]\label{thm:upper248}
Every snake in $\Q_9$ has at most 248 edges. Thus
\[
\boxed{191\le a(9)\le248.}
\]
\end{theorem}
\begin{proof}
The one-vertex path is immediate. For $v=m+1\ge2$, the same incidence estimate used for Theorem~\ref{thm:upper251} gives $140v+4T\le36\,856$. Lemma~\ref{lem:aggregate-four} gives $4T\ge8m-4=8v-12$. Therefore
\[
148v\le36\,868<148\cdot250.
\]
It follows that $v\le249$, so $m\le248$. The lower bound is Theorem~\ref{thm:lower}.
\end{proof}



\section{Endpoint singleton squares and the bound 247}
\begin{theorem}[At least four singleton squares at an endpoint]\label{thm:endpoint-four}
If $x$ is an endpoint of a snake in $\Q_9$ having at least one edge, then $\alpha_S(x)\ge4$.
\end{theorem}
\begin{proof}
Translate $x$ to zero and let $a$ be its sole selected unit neighbor. Suppose, for contradiction, that its family $F$ of missing coordinate pairs has at most three members. Every member of $F$ is a two-element subset of the eight coordinates other than $a$. For distinct $i,j\ne a$, the pair vertex $ij$ is selected if and only if $\{i,j\}\notin F$.

Call a coordinate isolated if it is different from $a$ and occurs in no member of $F$. The members of $F$ cover at most six coordinates, so there are at least two isolated coordinates. If $u,v$ are distinct isolated coordinates, then $uv$ is selected and has at most one selected neighbor, namely $auv$. Its unit neighbors are absent. For every other coordinate $k\notin\{a,u,v\}$, the three vertices $uv,uk,vk$ are selected, so $uvk$ is absent. Thus $uv$ is an endpoint.

There cannot be three isolated coordinates $u,v,w$: the three distinct vertices $0,uv,uw$ would all be endpoints. Consequently there are exactly two isolated coordinates, say $u,v$. The other six coordinates, apart from $a$, must be covered by the at most three missing pairs. They are therefore covered by exactly three pairwise disjoint pairs. In particular, each coordinate $t\notin\{a,u,v\}$ has a unique partner $j$, with $\{t,j\}\in F$ and $j\notin\{a,u,v,t\}$.

The selected vertices $0$ and $uv$ already account for the only two possible endpoints. Fix any $t\notin\{a,u,v\}$ and let $j$ be its partner. The selected pair $tu$ differs from both endpoints and is therefore internal. Its only possible selected neighbors are $atu$ and $tju$: for a third coordinate $k\notin\{a,t,u,j\}$, all three pair vertices $tu,tk,uk$ are selected, making $tuk$ absent. Both $atu,tju$ must therefore be selected. Applying the same argument to $tv$ forces $atv$ to be selected.

It follows that $at$ is absent for every $t\notin\{a,u,v\}$, since otherwise it would have the three distinct selected neighbors $a,atu,atv$. Choose two distinct coordinates $r,s$ outside $\{a,u,v\}$. The preceding argument also gives the selected triples $aru,asu$, so $au$ is absent because its three neighbors $a,aru,asu$ would be selected. Likewise $av$ is absent, using $arv,asv$.

All pair vertices $at$ with $t\ne a$ are now absent. Thus the selected vertex $a$ has only the selected neighbor zero and is itself an endpoint. The vertices $0,uv,a$ are distinct, giving a third endpoint and the required contradiction.
\end{proof}

\begin{lemma}[Aggregate count with endpoint contributions]\label{lem:aggregate-endpoints}
For a snake with $m\ge1$ edges and total singleton count $T$, one has $4m+10\le2T$.
\end{lemma}
\begin{proof}
For $m=1$, both endpoints have count at least four, so $2T\ge16\ge14$. Let $m\ge2$. Each of the first and last adjacent pairs has singleton counts summing to at least five, by Theorem~\ref{thm:endpoint-four} and the universal positive count. Each intervening adjacent internal pair has sum at least four by Theorems~\ref{thm:no-adjacent-singletons} and~\ref{thm:no-adjacent-one-two}. Therefore
\[
2T-\alpha_0-\alpha_m\ge4(m-2)+10=4m+2.
\]
Adding the endpoint contributions $\alpha_0+\alpha_m\ge8$ proves the result.
\end{proof}

\begin{theorem}[Verified upper bound 247]\label{thm:upper247}
Every snake in $\Q_9$ has at most 247 edges. Hence
\[
\boxed{191\le a(9)\le247.}
\]
\end{theorem}
\begin{proof}
Again the one-vertex path is immediate. For $v=m+1\ge2$, combine $140v+4T\le36\,856$ with Lemma~\ref{lem:aggregate-endpoints}, which gives $4T\ge8m+20=8v+12$. Then
\[
148v\le36\,844<148\cdot249.
\]
Thus $v\le248$ and $m\le247$. Combine this with the explicit 191-edge witness.
\end{proof}



\section{A further local exclusion: adjacent counts one and three}
\begin{theorem}\label{thm:no-adjacent-one-three}
Adjacent internal vertices cannot have singleton-square counts one and three. Consequently, every internal neighbor of an internal count-one vertex has singleton-square count at least four.
\end{theorem}
\begin{proof}
Use the same normalization as for Theorem~\ref{thm:no-adjacent-one-two}: zero has selected unit neighbors $a,b$, the vertex $ac$ is selected, and $cd$ is the unique missing root pair. Let $F$ now contain the three missing coordinate pairs at $a$, all avoiding $a,c$. Write $W$ for the five coordinates outside $\{a,b,c,d\}$ and put $t_i=abi$.

Let $H=\{h:\{b,h\}\in F\}$, and let $G$ consist of the members of $F$ not containing $b$. Thus $|H|+|G|=3$. For $e\in G$, put $Q_e=e\setminus H$ when $e\cap H\ne\varnothing$, and $Q_e=\varnothing$ otherwise. Set
\[
B=H\cup\bigcup_{e\in G}Q_e.
\]
Since every $e$ has two elements, $|Q_e|\le1$ and hence $|B|\le3$. Call an index $i\in W\setminus B$ good. Such an index satisfies $\{b,i\}\notin F$, and no missing pair $\{i,j\}$ joins it to a coordinate $j\in H$. In particular $t_i$ is selected.

We first identify an endpoint group for each good index. Apart from $abci$, the only possible selected neighbors of $t_i$ have the form $abij$, with $j\notin\{a,b,c,i\}$. A selected $abij$ requires $\{b,j\}\in F$ or $\{i,j\}\in F$; otherwise it would have the three selected neighbors $t_i,t_j,aij$.

Suppose $j\in W$ and $\{b,j\}\notin F$. Then $t_j$ is selected, so $aij$ is absent and $abij$ forces $bij$ absent. The selected root pair $ij$ has only $aij,bij$ as possible selected neighbors: every other triple $ijk$ is excluded by the three selected pair vertices $ij,ik,jk$. Thus $ij$ would be isolated in a nontrivial path, which is impossible. If instead $j\in H\setminus\{d\}$, goodness gives $\{i,j\}\notin F$. The selected vertices $t_i,aij$ then force $bij$ absent, making the selected pair $ij$ an endpoint.

For the remaining coordinate $j=d$, if $\{b,d\}\notin F$, then $t_d$ is selected. A selected $abdi$ forces $\{d,i\}\in F$ and makes $di$ an endpoint by the earlier exceptional-row neighbor classification. If $\{b,d\}\in F$, goodness ensures $\{d,i\}\notin F$, so $adi$ is selected. Unless an endpoint has already occurred among the pairs $hi$ with $h\in H\setminus\{d\}$, only $abci,abdi$ remain as possible selected neighbors of $t_i$. The eight-cycle argument from the preceding adjacent-count proof then forces an endpoint among $t_i,ci,di$.

Therefore every good index $i$ has an endpoint in the selected-vertex group consisting of $t_i$, all pairs $hi$ for $h\in H\setminus\{d\}$, the pair $di$ when $\{d,i\}\in F$ or $\{b,d\}\in F$, and the pair $ci$ when $\{b,d\}\in F$. Groups for different good indices are disjoint: each vertex in a group contains that group's index and no other good index. Thus there cannot be three good indices.

There are at least two good indices because $|W|=5$ and $|B|\le3$. If there are only two, say $u,v$, equality forces $|B|=3$ and $B\subseteq W$. Every member of $G$ must meet $H$, since otherwise its empty contribution $Q_e$ would make $|B|\le|H|+|G|-1=2$. Every edge of $G$ therefore lies in $B$, and every other edge of $F$ lies in $\{b\}\cup B$. In particular $d,u,v$ occur in no member of $F$.

The two endpoint groups now reduce to $\{t_u\}\cup\{hu:h\in H\}$ and its $v$ analogue. A third, disjoint selected group is $\{cu,adu\}$. To see that it contains an endpoint, note that $adu$ has only $du,acdu$ as possible selected neighbors. Indeed, for any other eligible coordinate $j$, the isolation of $d,u$ in $F$ makes $adj,auj$ selected, excluding the quad $aduj$ by its three selected neighbors. If $adu$ is internal, it forces $acdu$ selected. Then $cdu$ is absent because its neighbors $cu,du,acdu$ are selected. If $cu$ is also internal, its remaining neighbor possibilities force $acu$ selected. But $acu$ would have the three selected neighbors $ac,cu,acdu$. Thus $cu$ or $adu$ is an endpoint, giving the impossible third endpoint.

This proves the exclusion. The final assertion combines it with the previous exclusions of adjacent counts $(1,1)$ and $(1,2)$, and positivity of the singleton count.
\end{proof}



\section{Three consecutive count-two vertices and the bound 245}

\begin{theorem}[No three consecutive twos]\label{thm:no-three-twos}
There are no three consecutive internal vertices of a snake in $\Q_9$ whose singleton-square counts are all two.
\end{theorem}

\paragraph{Proof outline by finite geometric classification.}
The theorem is established by a finite geometric classification, formally checked in Lean. The following proof outline gives the normalization, classification, and principal contradictions. It does not reproduce every branch of the formal derivation. All coordinate labels in the classification are arbitrary distinct coordinates, and no solver infeasibility result is assumed.

Translate the first of the three vertices to zero and write the five consecutive selected vertices as
\[
c,\quad 0,\quad a,\quad ab,\quad abd.
\]
Here concatenation denotes exclusive-or of coordinate units. Inducedness implies that $a,b,c$ are distinct and $d$ differs from $a,b$. Thus either $d=c$ or all four coordinates are distinct. At each of the three internal centers, count two means that precisely two eligible unordered coordinate pairs have absent opposite corners. Each missing family has exactly two distinct members.

Three facts drive the case analysis. A selected vertex has at most two selected neighbors; at most two selected vertices have at most one selected neighbor; and a selected cycle is impossible. In particular, if the selected middle vertex $aij$ has both outer neighbors $ij,abij$, every other neighbor of $aij$ is absent. This saturation rule sharply limits ordinary extensions. Endpoint demands are counted only after proving their vertex sets disjoint, using coordinate bits.

If $d=c$, there are six remaining coordinates. Counting possible extensions of the selected rows $bk$, allowing at most two endpoint rows, forces the two outer missing families to be the same matching of two disjoint pairs. The two unused coordinates $u,v$ give endpoints $bu,bv$. The middle missing family must equal that matching: otherwise a selected middle vertex associated with a matching edge has at most one possible neighbor and supplies a third endpoint. The selected vertices $acu,acv$ are consequently internal and force $cu,cv$, giving the selected vertex $c$ the three neighbors $0,cu,cv$.

Now let $a,b,c,d$ be distinct, and let $R$ contain the other five coordinates. Denote the outer missing families by $F_0,F_2$, with the second expressed relative to the center $ab$. The common absent-row set is
\[
H=\{k\in R:\{b,k\}\in F_0\}
 =\{k\in R:\{a,k\}\in F_2\}.
\]
Put $h=|H|$, $\delta_0=[\{b,d\}\in F_0]$, and $\delta_2=[\{a,c\}\in F_2]$. Removing these specified edges leaves budgets $2-h-\delta_0$ and $2-h-\delta_2$. Every selected row $bk$ can continue only toward $bck,bdk$, or an ordinary extension $bkl$. Counting which missing edges can permit these continuations gives an exhaustive classification. The case $(\delta_0,\delta_2)=(0,0)$ is impossible. If exactly one indicator is one, reflection reduces to
\[
F_0=\{bd,ij\},\qquad F_2=\{ij,ck\},
\]
where $i,j,k\in R$ are distinct. The two remaining row endpoints force the middle family to be $\{ck,ij\}$, after which a fresh middle vertex has at most one neighbor, a third endpoint.

If both indicators are one, a common row hole gives $F_0=\{bd,be\}$ and $F_2=\{ac,ae\}$. Otherwise write
\[
F_0=\{bd,E_0\},\qquad F_2=\{ac,E_2\},
\]
where $E_0\subseteq\{d\}\cup R$ and $E_2\subseteq\{c\}\cup R$. The remaining branches are as follows.
\begin{itemize}
\item In the common-hole case, isolated middle rows and missing ordinary corners require at least three disjoint endpoint groups.
\item If $E_0=de$ and $E_2=cf$ are both stars, endpoint counting leaves two isolated groups; the resulting internal vertices force a vertex with three selected neighbors. This includes $e=f$.
\item If one remainder is a star and the other is an ordinary pair, every middle-isolated row and every ordinary middle missing edge supplies a disjoint endpoint resource. A two-edge family on five coordinates requires at least three such resources. Reflection handles the other orientation.
\item If $E_0=E_2=ij$ is ordinary, the three coordinates outside $i,j$ first force $ij$ to be middle-missing. Every possible second middle edge then forces three endpoints or a degree-three selected vertex.
\item Distinct ordinary remainders are treated by the endpoint count below.
\end{itemize}

\subsection*{The endpoint count for distinct ordinary remainders}
Consider the normalized window $c,0,a,ab,abd$, where $a,b,c,d$ are distinct,
and let $R$ be the other five coordinates. In this branch the outer missing
families are
\[
 F_0=\{\{b,d\},A\},\qquad F_2=\{\{a,c\},B\},
 \qquad A,B\in\binom{R}{2},\quad A\ne B.
\]
Here $F_2$ is expressed relative to the center $ab$. Let
$F\subseteq\binom{\{c,d\}\cup R}{2}$ be the two missing pairs at the middle
vertex $a$, so that $|F|=2$ and
\[
 aij\text{ is selected}\quad\Longleftrightarrow\quad\{i,j\}\notin F
 \qquad(i,j\notin\{a,b\},\ i\ne j).
\]
As throughout the argument, a string of coordinate labels denotes their
exclusive-or. For a two-element set $S=\{i,j\}$ we also write $[S]=ij$ and
$a[S]=aij$.

A coordinate $k\in R$ is \emph{middle-isolated} if it belongs to no member
of $F$. For $s\in\{c,d\}$ define the condition
\[
 C_s(k):\quad
 \text{for every }r\in R,\quad
 \{k,r\}\in\{A,B\}\ \Longrightarrow\ \{s,r\}\notin F.
\]
Call $k$ a \emph{good row} if it is middle-isolated and
\[
 \bigl(C_c(k)\wedge C_d(k)\bigr)
 \quad\vee\quad
 \bigl(\{c,d\}\notin F\ \wedge\ (C_c(k)\vee C_d(k))\bigr).
\]
Write $I$ for the set of good rows, and set
\[
 L=\{S\in F:S\subseteq R,\ S\cap A=\varnothing\},\qquad
 Q=\{S\in F:S\subseteq R,\ S\cap B=\varnothing\}.
\]

For $(H,K)=(A,B)$ or $(B,A)$, say that the outer hole $H$ is
\emph{eligible} when all three conditions hold:
\begin{enumerate}
\item $H\notin F$;
\item $\{c,s\},\{d,s\}\notin F$ for every $s\in H$;
\item for every $r\in R\setminus H$, some $s\in H$ satisfies
      $\{s,r\}\notin F$ and $\{s,r\}\ne K$.
\end{enumerate}
Let $\varepsilon_A,\varepsilon_B$ be the indicators of these respective
eligibility conditions. The endpoint demand is the explicit integer
\begin{equation}\label{eq:ordinary-endpoint-demand}
 D=|I|+|L|+|Q|+\varepsilon_A+\varepsilon_B.
\end{equation}

\paragraph{Why these quantities count endpoints.}
Every $bk$, with $k\in R$, is selected. For a good row $k$, the selected
set
\[
 \Gamma_k=\{bk,ack,adk\}
\]
contains an endpoint. To see the mechanism, an ordinary extension $bkr$
is excluded: $bk,br$ are selected, and at least one of $kr,abkr$ is selected
because $A\ne B$. Thus an internal $bk$ forces both $bck$ and $bdk$.
The vertex $ck$ is then absent, since its three neighbors $c,ack,bck$
would otherwise be selected. Likewise $abdk$ is absent, using its three
neighbors $abd,adk,bdk$.

The condition $C_c(k)$ leaves only $abck,acdk$ as possible selected
neighbors of $ack$; $C_d(k)$ leaves only $dk,acdk$ for $adk$. For an
ordinary extension, this follows either from saturation of a selected
middle triple with both outer corners present, or from three selected
middle triples when its pair is $A$ or $B$. If both conditions hold and
none of $\Gamma_k$ is an endpoint, the forced continuations form the
selected eight-cycle
\[
 bk,\ bck,\ abck,\ ack,\ acdk,\ adk,\ dk,\ bdk,\ bk.
\]
If instead $\{c,d\}\notin F$, the vertex $acdk$ is absent because
$acd,ack,adk$ are selected. Either clear direction then already gives
an endpoint in $\Gamma_k$.

Each $S=\{i,j\}\in L$ gives the endpoint $[S]=ij$. It is selected, while
$aij$ is absent. Its unit neighbors are absent; the extension $bij$ is
excluded by $bi,bj,ij$, and $dij$ by $di,dj,ij$. Every other ordinary
extension $ijr$ is excluded by $ij,ir,jr$, since $S$ is disjoint from $A$.
Only $cij$ can remain. Reflection gives the endpoint $ab[S]$ for each
$S\in Q$.

An eligible hole $H=\{i,j\}$ gives the selected endpoint $a[H]$.
Exactly one of its outer neighbors $[H],ab[H]$ is selected. The second
eligibility condition excludes its $c$- and $d$-extensions by the
three-neighbor rule. For each other coordinate $r$, the third condition
provides a selected middle triple $asr$ with both outer corners present;
its saturation excludes the extension $aijr$. There is therefore at
most one selected neighbor of $a[H]$.

\paragraph{Why the endpoint counts do not overlap.}
Different row groups $\Gamma_k$ are disjoint: every member contains
exactly the ordinary coordinate $k$. Moreover, $k\in I$ lies in no
member of $F$, so that group cannot meet either corner counted by $L$
or $Q$. Root corners $[S]$, translated corners $ab[S]$, and hole
vertices $a[H]$ have respectively the $(a,b)$ bit patterns $(0,0)$,
$(1,1)$, and $(1,0)$. These distinguish the three kinds. Within each
kind, distinct pairs give distinct vertices, including the two holes
because $A\ne B$. Finally a row vertex $bk$ has $a$-bit zero, whereas
$ack,adk$ have respectively $c$- or $d$-bit one; neither can equal a
hole vertex, whose $c,d$ bits are both zero. Thus choosing one endpoint
from each good-row group and adding the other counted endpoints gives
$D$ distinct endpoints. A path has at most two, so
\begin{equation}\label{eq:ordinary-demand-upper}
 D\le2.
\end{equation}

\subsection*{The final finite classification: scope of the outline}
The remaining finite argument has two parts. If $A$ and $B$ are disjoint,
two-edge counting gives $D\ge3$, contradicting
\eqref{eq:ordinary-demand-upper}. Its cases are indexed by the number
of members of $F$ wholly contained in $R$, which is zero, one, or two.
The isolated rows, clean corners, and eligible-hole terms in
\eqref{eq:ordinary-endpoint-demand} supply the required three endpoints.
The intermediate counting inequalities of these cases are not expanded
in this outline.

Otherwise $A=\{e,x\}$ and $B=\{e,y\}$ for distinct $e,x,y\in R$.
Let $u,w$ be the remaining two coordinates. Every middle missing edge
is one of the 21 pairs from $\{c,d,e,x,y,u,w\}$. Substituting these
possibilities in the explicit definition of $D$, and retaining two
distinct edges with $D\le2$, leaves exactly the following twenty
labelled families:
\[
\begin{gathered}
 \{cd,ce\},\ \{cd,de\},\ \{ce,de\},\ \{cx,dy\},\ \{cy,dx\},\\
 \{cx,ey\},\ \{dx,ey\},\ \{cy,ex\},\ \{dy,ex\},\\
 \{cr,xy\},\ \{dr,xy\}\quad(r=e,u,w),\\
 \{ex,ey\},\ \{ex,xy\},\ \{ey,xy\},\ \{eu,xy\},\ \{ew,xy\}.
\end{gathered}
\]
There are $21^2=441$ ordered choices, of which 21 repeat the same edge. The remaining 420 ordered choices represent 210 unordered two-edge families. The twenty displayed families are the unordered survivors of the endpoint-demand inequality. This classification does not assume a symmetry quotient. Each
family has a geometric exclusion by the degree-two and two-endpoint
conditions. The central and crossed-star families force a third neighbor
or a third endpoint after their clean row groups are accounted for.
For $\{dr,xy\}$, the required internal continuations force $acxy$ to
have the selected neighbors $acx,acy,cxy$; reflection treats the
$\{cr,xy\}$ family. The same mechanism covers $\{ex,xy\}$ and
$\{ey,xy\}$. For $\{ek,xy\}$ with $k=u$ or $w$, the other clean row
requires one endpoint. If $aexy$ is absent, $aex,aey$ provide two more;
if it is selected, saturation instead makes $xy,abxy$ two more.
The residual deductions complete the formal classification outlined here.


\begin{lemma}[Aggregate bound with forbidden runs]\label{lem:run-potential}
Let $v\ge2$ be the number of vertices of a snake and let $T$ be its total singleton count. Then
\[
3T\ge7v+5.
\]
\end{lemma}
\begin{proof}
Write the counts as $\alpha_0,\ldots,\alpha_{v-1}$. They are positive, both endpoint counts are at least four, and every neighbor of a count-one vertex has count at least four. Theorem~\ref{thm:no-three-twos} excludes every run $2,2,2$; a run involving an endpoint is also excluded by the endpoint bound.

Define a nonnegative potential
\[
P_i=\begin{cases}
4,&\alpha_i=1,\\
2,&\alpha_i=2,\ i>0,\ \alpha_{i-1}=2,\\
1,&\alpha_i=2\text{ in every other case},\\
0,&\alpha_i\ge3.
\end{cases}
\]
For each $i<v-1$ the local restrictions give
\[
7+P_i\le3\alpha_{i+1}+P_{i+1}.
\]
Indeed, if $\alpha_{i+1}=1$, then $\alpha_i\ge4$ and both sides equal seven. If $\alpha_{i+1}=2$ and $\alpha_i=2$, the previous count cannot also be two, so $P_i=1$, $P_{i+1}=2$ and both sides equal eight. If $\alpha_{i+1}=2$ but $\alpha_i\ne2$, then $\alpha_i\ge3$, giving seven on both sides. Finally, if $\alpha_{i+1}\ge3$, either $P_i\le2$, or $\alpha_i=1$ forces $\alpha_{i+1}\ge4$; the inequality follows in either case.

The endpoint bounds give $P_0=P_{v-1}=0$ and $3\alpha_0\ge12$. Summing the $v-1$ inequalities therefore gives
\[
3T\ge12+7(v-1)=7v+5.
\]
This argument does not assume that every three-term window sums to at least seven: the abstract sequence $1,4,1$ satisfies the numerical restrictions used in this lemma and shows why that inference would be invalid. This observation does not assert that the sequence occurs along an actual snake.
\end{proof}

\begin{theorem}[Verified upper bound 245]\label{thm:upper245}
Every snake in $\Q_9$ has at most 245 edges. Hence
\[
\boxed{191\le a(9)\le245.}
\]
\end{theorem}
\begin{proof}
The zero-edge case is immediate. Otherwise put $v=m+1\ge2$. The established incidence identity, nonnegativity of $E$, and $W\le v-2$ give
\[
4T+144v\le36864+4(v-2),\qquad T+35v\le9214.
\]
Combining this with Lemma~\ref{lem:run-potential} gives
\[
112v+5\le27642,\qquad112v\le27637.
\]
Since $112\cdot247=27664$, one has $v\le246$ and $m\le245$. The lower bound is the unchanged explicit 191-edge witness.
\end{proof}



\section{Cycles as an alternative construction route}
\begin{theorem}[Deleting a cycle vertex]
Deleting one vertex from an induced cycle of $C$ vertices leaves a snake of length $C-2$.
\end{theorem}
\begin{proof}
Order the cycle as $v_0,\ldots,v_{C-1},v_0$ and delete $v_{C-1}$. Among the remaining vertices, the only surviving cycle edges join consecutive indices. The closing edge involved the deleted vertex; the induced-cycle hypothesis forbids all other edges. The resulting path has $C-1$ vertices and $C-2$ edges.
\end{proof}
\begin{corollary}
An induced cycle with 194 vertices in $\Q_9$ would disprove the upper bound 191.
\end{corollary}
\begin{proof}
Deletion gives a 192-edge snake.
\end{proof}
In particular, a 192-vertex coil gives only a 190-edge snake through this deletion operation. Edge and vertex counts cannot be interchanged.

\section{The two-layer construction and its obstruction}
Write a vertex of $\Q_9$ as $(v,b)$ with $v\in\Q_8$ and $b\in\{0,1\}$. Vertices in the same layer retain their original adjacency. Across the layers,
\[
(u,0)\sim(v,1)\quad\Longleftrightarrow\quad u=v.
\]
The same description defines the product of an arbitrary graph with a binary edge. Taylor discusses such lifts in \cite{taylor}.

\begin{theorem}[Exact joining condition]\label{thm:join}
Let $P=(p_0,\ldots,p_r)$ and $R=(q_0,\ldots,q_s)$ be induced paths in a graph. The ordered sequence
\[
(p_0,0),\ldots,(p_r,0),(q_0,1),\ldots,(q_s,1)
\]
is an induced path precisely when
\[
p_i=q_j\quad\Longleftrightarrow\quad i=r\ \hbox{and}\ j=0.
\]
When valid, its length is $r+s+1$.
\end{theorem}
\begin{proof}
The layer labels ensure that vertices in different layers are distinct. Each individual path retains its distinctness and induced adjacency conditions. Every cross-layer edge is caused by equality of the base vertices. In the displayed order there is exactly one consecutive pair across layers, at $(p_r,0),(q_0,1)$. Therefore the induced condition requires exactly the stated equality and forbids all other equalities. These conditions also suffice, since all possible pairs have been covered. The vertex count is $(r+1)+(s+1)$, hence the edge count is $r+s+1$.
\end{proof}

\begin{corollary}[Single-crossing length constraints]
Assume the known bound $a(8)=98$. A single-crossing snake in $\Q_9$ has length at most 197. Length 196 requires layer lengths $(98,97)$ or $(97,98)$; length 197 requires $(98,98)$.
\end{corollary}
\begin{proof}
The two layer segments are eight-dimensional snakes. Thus $r,s\le98$ and $L=r+s+1$. For $L=196$, $r+s=195$ forces the two possibilities; for $L=197$, both upper bounds must be attained.
\end{proof}
\begin{theorem}[Integer realization of the layer construction]
The map $J(u,b)=u+256b$, for $0\le u<256$ and $b\in\{0,1\}$, is injective and preserves adjacency in both directions between $\Q_8\mathbin{\square}K_2$ and the integer encoding of $\Q_9$. Consequently, compatible layer paths of lengths 98 and 97 give a 196-edge snake under the same definition as the verified witness. Compatible paths of lengths 98 and 98 give a 197-edge snake and would refute the conjectured upper bound 196.
\end{theorem}
\begin{proof}
The intervals $[0,255]$ and $[256,511]$ are disjoint, so equality of images forces equality of both components. Flipping a coordinate $k<8$ preserves the layer, because
\[
(u+256)\mathbin{\oplus}2^k=(u\mathbin{\oplus}2^k)+256,
\qquad u\mathbin{\oplus}2^k<256.
\]
Flipping coordinate 8 exchanges $u$ and $u+256$. These are all nine possible flips. Thus adjacency within a layer is exactly eight-dimensional adjacency, while adjacency between layers requires equal base vertices. Compose this map with Theorem~\ref{thm:join}. The lengths become $98+97+1=196$ and $98+98+1=197$. Conversely, every unwanted equality between base vertices creates a nonconsecutive cross-layer edge, so the compatibility condition remains necessary after applying $J$.
\end{proof}
This argument covers snakes that cross a chosen coordinate exactly once. A general snake can cross every coordinate repeatedly. Its intersection with a layer can then be a disjoint union of paths, so bounding that entire intersection by a single-snake length would be invalid.

The length-98 seed in Appendix~\ref{app:witness} is taken from \cite{records} and independently checked against all pairwise conditions. Two tests were performed. First, truncate either endpoint of an isometric copy to length 97, choose both joining endpoints, and enumerate all $8!$ coordinate permutations. Translation is forced by identifying the chosen endpoints. The $2\cdot2\cdot2\cdot8!=322\,560$ placements, counted with possible duplication, have minimum overlap 15. Theorem~\ref{thm:join} requires overlap one. Second, a constraint solver searched for an arbitrary length-97 partner avoiding every seed vertex except the chosen endpoint. Both endpoint instances were infeasible. These restricted computational exclusions have not been independently certified; other length-98 seeds and multiple-crossing constructions remain possible.

\section{Independent hypothesis from the first eight terms}
\begin{center}
\begin{tabular}{rrrr}
\toprule
$n$ & $a(n)$ & $3\cdot2^{n-3}$ & correction\\
\midrule
3&4&3&1\\4&7&6&1\\5&13&12&1\\
6&26&24&2\\7&50&48&2\\8&98&96&2\\
\bottomrule
\end{tabular}
\end{center}
The baseline coefficient is also recovered from successive differences within the two triples:
\[
\frac{7-4}{2^4-2^3}=\frac{13-7}{2^5-2^4}
=\frac{50-26}{2^7-2^6}=\frac{98-50}{2^8-2^7}=\frac38.
\]
Continuing the correction as $1,1,1,2,2,2,4,4,4,\ldots$ gives
\begin{equation}\label{eq:model}
\widehat a(n)=3\cdot2^{n-3}+2^{\lfloor(n-3)/3\rfloor},\qquad n\ge3.
\end{equation}
The terms at dimensions one and two are assigned separately as 1 and 2.

\begin{conjecture}\label{conj:196}
The maximum snake length in $\Q_9$ is $196$ edges.
\end{conjecture}
This numerical conjecture is motivated only by the displayed extrapolation; it is not used in any proof below. Equation~\eqref{eq:model} predicts $3\cdot64+4=196$.

\begin{proposition}[Model recurrence]
For $n\ge3$, the model satisfies
\[
\widehat a(n+3)+6\cdot2^{\lfloor(n-3)/3\rfloor}=8\widehat a(n).
\]
\end{proposition}
\begin{proof}
Put $t=n-3$ and $c=2^{\lfloor t/3\rfloor}$. Since $\lfloor(t+3)/3\rfloor=\lfloor t/3\rfloor+1$, the left side is
$24\cdot2^t+2c+6c=8(3\cdot2^t+c)$.
\end{proof}
At $n=6$ the recurrence predicts $8\cdot26-12=196$. This is the same evidence as the correction table, not an independent argument.

\begin{proposition}[The known terms do not determine the ninth prediction]
The model
\[
\widetilde a(n)=3\cdot2^{n-3}+1+\left\lfloor\frac{n-3}{3}\right\rfloor
\qquad(n\ge3)
\]
agrees with every known term, after assigning the first two terms separately, and gives $\widetilde a(9)=195\ne196$.
\end{proposition}
\begin{proof}
For $n=3,4,5$, both corrections equal 1. For $n=6,7,8$, both equal 2. At $n=9$, the corrections are respectively 3 and 4. Direct substitution proves the claim.
\end{proof}
Only two correction blocks have been observed. No construction currently explains why either proposed correction should describe the actual extremal function.

\subsection{A restriction on one-singleton configurations}
A possible route to a stronger aggregate bound is to show that vertices with only one singleton square must lie near an endpoint. The following result narrows the local configuration without establishing that claim.

\begin{theorem}[The two continuations use different missing coordinates]\label{thm:split-arms}
Let $S$ be an induced path vertex set in $\Q_9$ containing zero. Suppose its selected unit vectors are exactly $e_a,e_b$, where $a\ne b$, and put $U=\{0,\ldots,8\}\setminus\{a,b\}$. Let $u,v\in U$ be distinct. Assume every pair vertex $e_s+e_t$ with distinct $s,t\in U$, other than possibly $e_u+e_v$, is selected, and every selected pair vertex indexed by $U$ is internal in the path. Suppose also that $e_a+e_i$ and $e_b+e_j$ are selected for some $i,j\in U$. Then
\[
(i,j)=(u,v)\quad\text{or}\quad(i,j)=(v,u).
\]
\end{theorem}
\begin{proof}
First suppose that $i\notin\{u,v\}$. Choose two distinct coordinates $k$ outside $\{a,b,i,u,v\}$; at least four are available. For either choice, $e_i+e_k$ is selected and internal. Its only possible selected neighbors are $e_a+e_i+e_k$ and $e_b+e_i+e_k$. Indeed, deleting $i$ or $k$ gives an absent unit vector. Adding any other coordinate $t\in U$ gives a triple with three selected pair neighbors: none of its three pairs can be the exceptional pair $\{u,v\}$, since both $i$ and $k$ lie outside that pair. Such a triple is forbidden by degree at most two.

Both possible neighbors of the internal vertex $e_i+e_k$ must therefore be selected. The two choices of $k$ give two distinct neighbors of $e_a+e_i$, in addition to its selected neighbor $e_a$, a contradiction. Hence $i\in\{u,v\}$. Interchanging $a,b$ proves $j\in\{u,v\}$.

Suppose next that $i=j=u$. Choose three distinct $k\in U\setminus\{u,v\}$. Each selected internal vertex $e_u+e_k$ has possible selected neighbors only among
\[
e_a+e_u+e_k,\qquad e_b+e_u+e_k,\qquad e_u+e_v+e_k.
\]
All other triples have three selected pair neighbors, while unit neighbors are absent. At least one of the first two displayed vertices must therefore be selected, since an internal vertex requires two selected neighbors. However, $e_a+e_u$ already has selected neighbor $e_a$, and can have at most one of the three vertices $e_a+e_u+e_k$ as another selected neighbor. The same argument applies to $e_b+e_u$. Thus at most two of the three choices of $k$ can be covered, a contradiction. Interchanging $u,v$ excludes $i=j=v$. The two remaining ordered pairs are exactly those asserted.
\end{proof}

The theorem has an independent machine check under its stated geometric premises. If zero has exactly one singleton square and all endpoints have Hamming distance at least three from zero, the missing square supplies $u,v$, and every selected pair vertex is internal. The root and its two unit neighbors are internal as well; their second neighbors supply the two continuations. Thus the theorem identifies the remaining local arrangement. It does not prove that such a full path is impossible.

\begin{proposition}[Radius-three constraints do not suffice]\label{prop:radius-three}
There is an induced linear forest on 60 vertices of $\Q_9$, consisting of 14 pairwise separated paths, such that every selected vertex lies within Hamming distance three of zero, every selected vertex at distance at most two has degree two, and zero has exactly one singleton square.
\end{proposition}
\begin{proof}
Take the 14 ordered paths in Appendix~\ref{app:radius-three}. Direct checks show that every path is induced, their vertex sets are disjoint, and no edge joins different paths. Counting selected neighbors gives degree two at every selected vertex of weight at most two. Every listed vertex has weight at most three. Among the 36 coordinate squares through zero, exactly the square using coordinates two and three contains no other listed vertex. These finite checks have independent machine checks.
\end{proof}

The forest satisfies the local degree and acyclicity requirements but is disconnected. It is therefore not a counterexample to a statement about a single full snake. Its existence shows that these radius-three requirements alone cannot complete the proposed endpoint-proximity argument. Connections through higher-weight vertices, or another global restriction, remain necessary.

\subsection{Cycle constraints beyond four-cycles}
\begin{lemma}[A finite subset has a vertex of degree at most one]\label{lem:subset-low-degree}
Let $P=(p_0,\ldots,p_m)$ be an induced path and let $A$ be a nonempty subset of its vertex set. Some vertex of $A$ has at most one neighbor in $A$.
\end{lemma}
\begin{proof}
Choose the least index $i$ for which $p_i\in A$. Any neighbor $p_j\in A$ of $p_i$ satisfies $j\ge i$ by minimality and $|j-i|=1$ by the induced-path condition. Thus $j=i+1$. There is at most one such vertex.
\end{proof}
Consequently, a finite nonempty vertex set having two distinct neighbors within the set at every vertex cannot be contained in an induced path. This excludes cycles of arbitrary length and also applies to disconnected unions of cycles. The lemma and the validity of all 41 recorded cycle cuts in the radius-five experiment have been independently machine-checked.

\section{Computational investigations and their exact scope}
\subsection{Published bank audit and endpoint regrowth}
Independent Python checks verified all 131 published nine-dimensional examples, including their transition and vertex lists, every pairwise adjacency condition, and blocked endpoints. The total number of unordered pairs checked was $131\binom{192}{2}=2\,402\,016$. Canonical transition sequences, normalized by first coordinate occurrence and compared with their reversals, distinguish 131 classes under translation, coordinate permutation, and reversal. This audit makes no completeness claim.

Endpoint regrowth retained a prefix of each example in each orientation, then searched every legal continuation to 193 total vertices. A vertex is forbidden after it has been used, or when it is adjacent to a retained vertex preceding the current endpoint. Adding a vertex blocks all neighbors of the previous endpoint. The resulting depth-first search covers the permitted continuations. Node limits produce \textsc{Unknown} rather than a nonexistence result.
\begin{center}\begin{tabular}{rrrrr}
\toprule
Removed vertices & Retained vertices & Excluded cases & Unresolved cases & Nodes\\\midrule
32&160&262&0&488,450\\
40&152&262&0&16,561,030\\
48&144&209&53&66,968,134\\\bottomrule
\end{tabular}\end{center}
Each case is one of the 262 oriented prefixes obtained from 131 paths and their reversals. The final row used one million nodes per orientation. This table records the initial bounded search. Subsequent certificate work excludes 192-edge extensions for all 262 orientations while retaining the first 112 vertices, with a completed Lean verification receipt. The exclusion concerns these particular retained segments; it neither classifies all snakes nor prevents improvements that change the interior.

\subsection{Constraint models and vertex replacement}
The exact model selects cube vertices and directed cube-edge arcs. A single directed circuit through an auxiliary vertex encodes one connected path; deleting the auxiliary vertex gives its endpoints. Unselected cube vertices use self-loops. Every cube edge whose two endpoints are selected is required to appear in the path. Together these constraints enforce distinctness, connectivity, and absence of chords. Small feasible and infeasible thresholds in dimensions three through five were reproduced. The unrestricted 192-edge run returned \textsc{Unknown} after 300 seconds.

A broader local model permits both endpoints to change and imposes only
\[
|V(P)\setminus V(T)|\le k,\qquad |V(T)|=193,
\]
where $P$ is a recorded 191-edge snake. Thus changes may occur anywhere in the path. For $k=2$, all 131 models were infeasible, with about 241 seconds total solver time. In vertex-set terms, a 192-edge improvement would have to delete at least three vertices from each tested representative. A two-vertex deletion allows three added vertices, since the target has one extra vertex overall. This excludes a restricted neighborhood, without excluding a distant snake. The corresponding model with $k=8$ returned \textsc{Unknown} on the first three seeds within 30 seconds each. The implementations and exclusions have not been independently certified.

\subsection{Fixed-length walks with temporary defects}
To search directly for 196 edges, consider a walk $v_0,\ldots,v_L$ in $\Q_9$ with $\Ham(v_i,v_{i+1})=1$. Distinctness and the absence of chords are temporarily relaxed. Define
\[
D(v)=\bigl|\{(i,j):0\le i<j\le L,\ v_i=v_j\}\bigr|,
\qquad
C(v)=\bigl|\{(i,j):0\le i<j-1\le L-1,\ \Ham(v_i,v_j)=1\}\bigr|,
\]
and set $E(v)=4D(v)+C(v)$. The positive weight four is a search choice.

\begin{proposition}\label{prop:zero-defect}
An adjacent walk is an induced path if and only if $E(v)=0$.
\end{proposition}
\begin{proof}
Both counts are nonnegative. Therefore $E(v)=0$ is equivalent to $D(v)=C(v)=0$. The first equality gives distinct vertices. The second excludes all nonconsecutive adjacencies, while consecutive adjacency is already assumed. Conversely, an induced path has neither repeated vertices nor chords.
\end{proof}
This equivalence has been checked in Lean under the same path definition as the 191-edge witness. It supplies an acceptance criterion and does not establish that the search reaches zero.

Two types of move preserve consecutive adjacency. Exchanging two distinct transition directions translates the intervening vertex interval by their binary sum. Changing one transition direction translates the entire suffix by the sum of its old and new directions. Distances within the translated interval and its complement remain unchanged, so only pairs crossing its boundary affect the score. These moves permit interior changes and moving endpoints. Simulated annealing permits temporary increases in $E$.

Two 90-second runs at $L=196$ each reduced the best score from 10 to 7, after approximately 10.3 and 10.4 million proposals. Both saved walks have 197 distinct vertices and seven forbidden chords. A diagnostic run at $L=192$ retained its initial score of one after approximately 10.2 million proposals. No valid path longer than 191 edges was found. Independent pairwise checks confirmed the saved scores and valid segments. A separate implementation audit compared 104,939 evaluated move scores with direct pair counts. These computational checks do not constitute a proof of the search implementation.

Exact repair then released neighborhoods of all defects and both endpoints, requiring the remaining vertices to be selected but leaving their order unrestricted. For one 196-edge candidate, fixing 145 vertices gave infeasibility in the restricted model; fixing 105 vertices was unresolved after 90 seconds. For the other candidate, fixing 85 vertices was unresolved after 60 seconds. A 192-edge repair fixing 160 vertices was infeasible. A control fixing a known 191-edge path succeeded and passed an independent check of all 18,336 unordered vertex pairs. The infeasibility reports have no Lean certificate and exclude only their specified fixed vertex sets.

The score measures violations, rather than the number of edits needed to obtain a snake. Seven chords may require substantial changes to remove. The present results neither establish nor refute the proposed value 196.

\subsection{Diverse candidates and solver-guided interior repair}
A subsequent comparison retained either one walk or up to 32 distinct walks attaining the current best score. Translation, coordinate permutation, and reversal were removed when identifying duplicate candidates. Both variants used equal proposal budgets. With random seeds 41 and 113, the single-candidate runs reached scores 8 and 9 after ten million proposals, while the population runs reached 6 and 7. Extending each run to one hundred million proposals gave scores 6 and 6 for the single-candidate variant, and 6 and 7 for the population variant. Thus the initial advantage did not persist consistently. These eight runs used 440 million proposals in total and found no valid improvement over 191 edges.

Every saved best walk from the four longer runs has a 191-edge prefix equivalent to a member of the published bank. All its forbidden chords touch one of the final five vertices. This observation concerns the saved states and does not classify the intermediate search states. It indicates that reducing this score can return to a known path with a defective appended tail even when interior moves are permitted.

Two six-chord candidates were also subjected to adaptive repair. A minimum set of positions meeting every chord was removed first, leaving 192 or 194 initially fixed vertices. After an infeasible restricted model, the solver supplied a sufficient conflicting subset of the requirements that particular vertices remain selected. One policy released a vertex from that subset together with nearby positions along the walk; another released up to five scattered vertices from the subset. These subsets were not claimed to be minimum. Four runs ended unresolved, with 140, 97, 131, and 84 fixed vertices respectively. The preceding 17, 19, 20, and 22 infeasible rounds concern only their recorded fixed sets and have no Lean impossibility certificates. Different final constraints prevent interpreting these counts as a ranking of the policies. Independent checks verified the saved walks, the minimum initial chord covers, the successive releases, and a successful 191-edge control.

\subsection{A construction through cyclic components}
These results motivate a construction search that relaxes connectivity. Select 197 vertices, require exactly two selected vertices to have selected degree one, and require all other selected vertices to have selected degree two. A finite connected graph with degrees at most two is a path or a cycle. Since every path component here has two degree-one vertices and isolated vertices are excluded, the selected graph consists of one path and zero or more disjoint cycles.

For a cyclic component with vertex set $C$, add the constraint
\[
|C\cap V(P)|\le |C|-1.
\]
Lemma~\ref{lem:subset-low-degree} justifies this inequality for every induced path $P$. The cardinality inequality and its connection to the finite cycle checker have been checked in Lean. The search can then reconsider the entire vertex selection. Neither practical termination nor the existence of a 196-edge solution follows from this reduction.

A concrete starting configuration comes from two vertex-disjoint induced 26-cycles $A,B\subseteq\Q_6$. Put $A$ in each even-parity layer of $\Q_3\mathbin{\square}\Q_6$ and $B$ in each odd-parity layer. Adjacent layers have different parity, and a cross-layer edge requires equal projected vertices; disjointness therefore forbids all cross-layer edges. The product has eight induced 26-cycles and 208 vertices. Deleting eleven consecutive vertices from one cycle leaves one path on 15 vertices and seven intact cycles. Writing $P_r$ for a path on $r$ vertices and $C_r$ for a cycle on $r$ vertices, this gives
\[
P_{15}\sqcup 7C_{26},\qquad
|V|=15+7\cdot26=197,\qquad |E|=14+7\cdot26=196.
\]
The explicit 197-vertex set has passed Lean checks of its cardinality, distinctness, selected degrees, and partition; its seven cycle exclusions are also checked. An independent pairwise audit verifies the component sizes and the ordered path. These are facts about a disconnected graph and provide no improved snake lower bound. The construction uses cycles; the published 208-edge product of paths in \cite{orland} instead has 224 vertices.

Three initial 196-edge probes, with and without square cuts or a defective-walk hint, found no candidate within 20 seconds. The product hint supplied a feasible first-round configuration. Adding all seven cycle exclusions then gave an unresolved second round. Adding one cycle exclusion at a time produced a further candidate
\[
P_{41}\sqcup 6C_{26}
\]
after about 0.93 seconds, preserving all 197 selected vertices in number while replacing five vertices of the preceding candidate. Its path has 40 edges. The next round was unresolved after 20 seconds. Thus one component merger has been achieved in this experiment, while six cyclic components remain. The full solver implementation is not Lean-verified; its saved feasible sets and the returned paths have independent adjacency checks. Continuing from this configuration must retain the previously accumulated cycle exclusions.

\subsection{Frozen components and larger cycle configurations}
Further mergers retained five or four of the 26-cycles and replaced the remaining selected vertices by one path of 66 or 92 edges. Both accumulated cycle exclusions were preserved, and both path endpoints were free. Retaining the entire preceding 40-edge path gave infeasibility in the restricted model. Releasing that path, or allowing either of two pairs of cycles to change, remained unresolved after 60 seconds. The fixed-path diagnostic found no first extension from endpoint 120, and only the addition of vertex 2 from endpoint 0, with no further available continuation. These are restrictions of this intermediate configuration and provide no global bound.

If $F$ is the frozen vertex set and $R$ is the changing path, then $F\cap R=\varnothing$. A cycle exclusion for a set $C$ becomes
\[
|C\cap R|\le |C|-1-|C\cap F|.
\]
Indeed, $|C\cap(F\cup R)|=|C\cap F|+|C\cap R|$. The subtraction is over the integers: when $C\subseteq F$, the right-hand side is $-1$, correctly making the condition impossible. This equivalence and the fully frozen case have been checked in Lean. In particular, freezing a cycle that the latest cut asks the search to change is an invalid positive control; removing that pending cut recovers the known 40-edge path.

A larger starting configuration uses the following two cyclic orders in $\Q_5$, with vertices encoded by integers from 0 to 31:
\[
\begin{aligned}
A&=(0,1,17,25,27,31,15,13,12,28,20,22,18,2),\\
B&=(3,11,9,8,24,26,30,14,6,4,5,21,23,19).
\end{aligned}
\]
Both are induced 14-cycles, and their vertex sets are disjoint. Their lengths, distinctness, cyclic adjacency equivalences, and disjointness have been checked in Lean. Alternating $A$ and $B$ according to the parity of the sixteen cells of $\Q_4\mathbin{\square}\Q_5$ gives sixteen nonadjacent cycles and 224 vertices. As before, disjoint base sets exclude cross-cell edges. A separate size-28 base search in $\Q_6$ found two disjoint unions of two 14-cycles each, also giving 224 product vertices; it did not find 28-cycles. A size-54 pair search in $\Q_7$ was unresolved after 60 seconds.

The larger vertex count does not make a 197-vertex starting state automatic. Keeping whole 14-cycles and a nontrivial proper path from one further cycle leaves $14k+r$ vertices with $2\le r\le13$. Such a number cannot equal $197=14\cdot14+1$. This arithmetic obstruction has been checked in Lean and applies only to that trimming operation.

Joint cell replacement avoids the obstruction. Releasing cells $0,1,2$ and replacing their contents by a 14-edge path gave $P_{15}\sqcup13C_{14}$. Releasing either $0,1,2,3$ or $0,1,2,4$ gave $P_{29}\sqcup12C_{14}$. Each selected graph has 197 vertices and passed independent adjacency checks. Releasing the first five through eight cells, with target path lengths 42, 56, 70, and 84, remained unresolved after 60 seconds each. These disconnected constructions do not improve the snake lower bound.

\subsection{Connector costs and a finite cell model}
The four unused base vertices are
\[
U=\{7,10,16,29\}.
\]
Consider the restricted family in which each even cell permits $A\cup U$ and each odd cell permits $B\cup U$. Each cell has 18 available vertices. Since $(A\cup U)\cap(B\cup U)=U$, only the four shared vertices can carry an edge between adjacent cells. Their numbers of neighbors on the two base cycles are
\[
\begin{array}{c|rrrr}
u&7&10&16&29\\\hline
|N(u)\cap A|&1&1&4&4\\
|N(u)\cap B|&4&4&1&1
\end{array}
\]
and these finite counts have been checked in Lean.

\begin{proposition}[A local connector cost]
Let a selected vertex of an induced path have four possible neighbors in a specified set $N$ and one selected neighbor outside $N$. At least three vertices of $N$ are absent from the path.
\end{proposition}
\begin{proof}
A selected vertex has at most two selected neighbors. The neighbor outside $N$ leaves room for at most one selected member of $N$. Since $|N|=4$, at least three members are absent.
\end{proof}
This statement has been checked in Lean. In the cell construction, a cross-cell edge at a connector with four base-cycle neighbors therefore requires omitting at least three of those neighbors. Different connector costs can share omitted vertices; adding them as independent losses would be unjustified.

An enumeration of all $2^{18}$ local subsets retained 4,973 degree-compatible states for each cell type. A state specifies its selected vertices, selected connectors and their internal degrees, number of nonconnector endpoints, and vertex count. Nonconnector vertices must have selected internal degree one or two, while connectors can have degree zero before external edges are added. At most two nonconnector endpoints are permitted. Internal cycles remain allowed. Every saved row passed an independent adjacency check; enumeration completeness is not Lean-certified.

The global construction model now permits all sixteen cells to change simultaneously. It imposes 197 selected vertices, selected degree one or two, and exactly two endpoints. Ten distinct cyclic components occur in each local table; their translations give 160 cycle cuts, all checked in Lean. A second model also requires one connected component. Both 196-edge searches returned \texttt{UNKNOWN} after 120 seconds. A length-maximizing version found a 101-edge induced path, independently checked by pairwise adjacency. This path is a control for the construction method and does not improve the known lower bound.

\paragraph{A restriction that must be removed.}
A separate enumeration, using edge lists and cycle detection, examined all $2^{18}$ subsets of each cell and reproduced both local tables exactly. Removing states containing a cycle leaves 4,142 states per type. If $e$ is the number of selected nonconnector vertices of internal degree one, the maximum selected size in this acyclic census is
\[
\begin{array}{c|rrr}
e&0&1&2\\\hline
\text{maximum selected size}&10&14&14.
\end{array}
\]
Consequently, every enumerated acyclic state satisfies $c\le10+4e$. For a single induced path in the restricted family, nonconnectors have no neighbors in other cells, so their internal endpoints are actual path endpoints. Therefore $\sum_{h=0}^{15}e_h\le2$, and the census gives
\[
\sum_{h=0}^{15}c_h\le160+4\sum_{h=0}^{15}e_h\le168.
\]
Lean checks the numerical profiles and this summation under the local census premise. The exhaustive correspondence between all geometric subsets and those profiles remains computational evidence rather than a Lean completeness theorem. This restriction concerns only the alternating $A\cup U$, $B\cup U$ family. It does not challenge the conjectured maximum 196 in $Q_9$.

The constructive consequence is to enlarge the allowed vertices within cells. A useful next search permits vertices from the opposite base cycle, which can create cross-cell connections beyond the four shared connectors. The local census explains why retaining the original four-connector restriction is an unsuitable route to a 196-edge construction.

\subsection{Expanded cells and alignment with known snakes}
The enlarged model permits all 512 cube vertices and places a budget on the number selected outside the original 288-vertex cell palette. To choose a starting configuration, each of the 131 published 191-edge witnesses was aligned with that palette. The computation examined all 126 choices of five internal coordinates, all 120 orders of those coordinates, 32 internal translations, and two outer translation parities. Permuting the four outer coordinates or translating them by a word of even weight leaves palette membership unchanged. Thus these 967,680 representatives suffice for this alignment objective. Within each internal coordinate set, identical palette patterns reduced the 7,680 representatives to 960 distinct patterns.

The computed minimum numbers of outside-palette vertices range from 29 to 38 across the bank. Their frequencies are
\[
\begin{array}{c|rrrrrr}
\text{outside vertices}&29&30&35&36&37&38\\\hline
\text{bank witnesses}&2&2&3&57&63&4.
\end{array}
\]
Every transformed witness passed a separate pairwise induced-path check. These minima concern the finite bank and the specified alignment objective; they do not give a lower bound on outside-palette use by arbitrary snakes.

The selected witness is bank entry 129, using zero-based indexing, with 29 outside-palette vertices. Its alignment sends input bits to output bits in the order $(1,5,4,2,7,0,3,6,8)$ and then translates by the word with decimal value 7. A coordinate permutation followed by translation preserves which pairs differ in exactly one bit. For this particular map, Lean checks an inverse, every coordinate-flip identity, injectivity, adjacency in both directions, and the induced transformation on snakes of any length. This certifies the relabeling; the bank-wide minimization remains computational.

The expanded model recovered the fixed 191-edge witness with an outside budget of 29 and rejected it with budget 28. A 196-edge search allowing up to 32 deleted seed vertices and 34 outside-palette vertices returned \texttt{UNKNOWN} after 120 seconds. Two 192-edge searches, each allowing 16 deletions, also remained unresolved after 120 seconds: one allowed 40 outside vertices and the other imposed no effective palette restriction.

Whole-cell replacement was tested separately. With the remaining vertices fixed to the selected witness, all sixteen one-cell searches and all 32 adjacent two-cell searches for exactly 192 edges returned \texttt{INFEASIBLE}. These are restricted solver results without Lean infeasibility certificates. A 191-edge control succeeded. Allowing an equal-length change in cells 0 and 2 replaced one vertex and produced a snake equivalent to bank entry 127 under translation, coordinate permutation and reversal. The rearrangement changes the starting configuration but adds no new symmetry class to the bank. The next constructive step is to vary the coordinate partition or replace connections across more cells simultaneously.

\subsection{Coordinate-face replacements and a two-arm construction}
Changing the internal coordinate partition gives 672 six-dimensional coordinate faces in $Q_9$, since $\binom96 2^3=672$. For the selected aligned witness, the search fixed every selected and unselected vertex outside each face and allowed its interior to change. All 672 searches for exactly 192 edges returned \texttt{INFEASIBLE}, as did the 672 searches for exactly 196 edges. A 191-edge control succeeded. These are restricted solver results without Lean infeasibility certificates.

The geometric meaning of this restriction can be stated exactly. For a nonempty set $D$ of changed vertices, let $J(D)$ contain the coordinates on which two members of $D$ differ.
\begin{proposition}[Support of a face replacement]
The set $D$ is contained in a six-dimensional coordinate face if and only if $|J(D)|\le6$. For selected vertex sets $S,T$, agreement outside a region $R$ is equivalent to $S\mathbin{\triangle}T\subseteq R$.
\end{proposition}
\begin{proof}
If $D$ lies in a face whose six free coordinates form $I$, all members agree outside $I$, so $J(D)\subseteq I$. Conversely, extend $J(D)$ to a six-element coordinate set $I$ and choose $a\in D$. Every member of $D$ agrees with $a$ outside $I$, so $D$ lies in the corresponding face. The second assertion follows directly from the definition of symmetric difference.
\end{proof}
These statements have been checked in Lean. They identify the geometry covered by the search; they do not certify its solver exclusions.

The analogous seven-dimensional sweep has $\binom97 2^2=144$ faces. At both target lengths, a two-second limit per face gave 142 exclusions and two unresolved cases. Longer runs excluded one further case at each target. The remaining face fixes coordinates 7 and 8 to zero, in the aligned coordinates, and contains the first 50 vertices of the seed. Its complement retains a single 142-vertex path $F$.

This remaining case admits a smaller constructive model. The only vertices of the face adjacent to the two endpoints of $F$ are 53 and 44. Removing all face vertices adjacent to the interior of $F$ leaves an induced subgraph of $Q_7$ with 78 vertices and 147 edges. Seek two vertex-disjoint induced paths $L,R$ in this subgraph, ending at 53 and 44 respectively, with no edges between the two paths. Joining them through $F$ gives an induced path with
\[
|V(L)|+142+|V(R)|
\]
vertices. The 196-edge target therefore requires $|V(L)|+|V(R)|=55$. Computationally, $F$ is replaced by one artificial edge between 53 and 44; every feasible output is expanded and independently checked in $Q_9$.

The two-arm model reproduced a 143-edge control. Its 196-edge feasibility run returned \texttt{INFEASIBLE}; the 192-edge run remained unresolved after 60 seconds. A 120-second length-maximizing run produced a verified 186-edge snake, using 45 vertices in the face and the fixed 142 vertices of $F$. This path is below the known record. It is not equivalent to a contiguous 186-edge subpath of any of the 131 stored witnesses under translation, coordinate permutation and reversal, according to the finite comparison. It supplies a different starting configuration for subsequent construction searches. The next search must allow changes to $F$ as well, since retaining the two-arm restriction already gave a solver exclusion at 196 edges. The compressed model and its infeasibility result are computational evidence, not a Lean-certified exclusion.

\subsection{Other searches and unresolved cases}
Four structured searches for a 194-vertex coil used quotient-cube constructions with extra generators of weights 2, 4, 6, and 8. All were unresolved at their time limits; a small control lifted to a valid 14-vertex coil. No coil improvement follows.

For the fixed eight-dimensional seed, searches for partner lengths 93, 92, and 91 at both endpoints were each unresolved after 30 seconds. A relaxation imposing only induced degree at most two admits 106 vertices in the permitted region, so that local degree condition alone is insufficient. Retaining 50 initial seed vertices in either orientation and exhaustively completing to length 98 recovered only the original completion, after 1,926,810 nodes in total. Retaining 42 vertices reached its twenty-million-node limit in both orientations. Neither result classifies all length-98 snakes.

\subsection{Joint search with both layer paths variable}
The fixed-seed obstruction does not exclude other pairs. A joint model therefore selects both projected paths, requires each to start at zero, and forbids every shared nonzero projected vertex. A circuit constraint with one auxiliary vertex enforces a single connected path in each layer; selected adjacent cube vertices must be consecutive in that path. The only cross-layer edge is then the intended join at zero.

The two first coordinate directions must differ, since a common first neighbor would be an additional shared vertex. A common coordinate permutation may therefore send those directions to coordinates zero and one. With and without this symmetry restriction, all 42 tested cases in dimensions two and three agreed with independent exhaustive enumeration.

\begin{proposition}[A checked starting pair]\label{prop:joint-seed}
There are induced paths of lengths 96 and 94 in $\Q_8$ whose vertex sets intersect only at their common initial vertex. Their two-layer join has 191 edges.
\end{proposition}
\begin{proof}
Use the two transition sequences in Appendix~\ref{app:joint-seed}, both starting at zero. They reconstruct 97 and 95 vertices. All $\binom{97}{2}=4\,656$ and $\binom{95}{2}=4\,465$ within-path pairs satisfy the distinctness and induced-adjacency conditions. Checking all $97\cdot95=9\,215$ cross pairs gives equality exactly for the two initial vertices. Reversing the first path and applying the two-layer joining criterion gives $96+94+1=191$ edges. These finite checks and the joining deduction have independent machine checks.
\end{proof}

The pair was extracted from a published nine-dimensional example: coordinate five changes only once, at edge index 96, when edge indices start at zero. Deleting that coordinate and applying a common translation and coordinate permutation gives the displayed pair. This is a decomposition of a published example, not a new record. Fixing both selected vertex sets to this pair gave a feasible solver control.

Three searches then left both paths variable. Each used four workers, random seed 13, and a 120-second limit:
\begin{center}
\begin{tabular}{rrrrl}
\toprule
First length & Second length & Joined target & Starting hint & Status\\
\midrule
98 & 97 & 196 & none & \texttt{UNKNOWN}\\
98 & 98 & 197 & none & \texttt{UNKNOWN}\\
96 & 95 & 192 & 96,94 pair & \texttt{UNKNOWN}\\
\bottomrule
\end{tabular}
\end{center}
The hint in the last run does not fix either path. These timeouts establish no existence or impossibility result. The model and its search implementation have not been independently certified; the explicit starting pair and the mathematical joining criterion have been checked separately.

\subsection{A radius-five exclusion and its verification boundary}
Fix the normalized radius-two selected set
\[
F=\{0,1,2,5,10\}\cup
\{2^a+2^b:2\le a<b\le8,\ (a,b)\ne(2,3)\},
\]
where integers encode binary vertices. In the ball of Hamming radius $r$ about zero, the model selects exactly $F$ at weights at most two. Every selected vertex strictly inside the ball has degree one or two; selected vertices of weight at most two have degree two. Boundary vertices have degree at most two. The number of interior vertices of degree one is bounded by a prescribed budget. Connectedness is not imposed. Cycle cuts exclude specified cyclic components.

At radius four, a forest exists with interior endpoint budget two. At radius five, the solver reports infeasibility for budgets two and three, and finds a forest with budget four. Allowing cycles restores feasibility at budget two. The feasible assignments have been checked independently. These findings distinguish an obstruction involving both endpoint count and cycles from the weaker radius-three conditions in Proposition~\ref{prop:radius-three}.

A Boolean encoding of the radius-five, budget-two model records the exact cycle cuts. Removing irrelevant clauses and replacing auxiliary counters by direct restrictions on triples of endpoint indicators gives a 567-clause core involving 252 Boolean variables. Its certificate consists of 398 derived clauses obtained by reverse unit propagation: assume the negation of the proposed clause, propagate the cited clauses, and obtain a contradiction. The resulting Boolean contradiction has been independently machine-checked through ordinary logical proof steps. This certifies the explicit formula; no geometric exclusion is asserted without the remaining interpretation proof.

The remaining verification boundary is explicit. The interpretation of each input clause as a necessary consequence of the geometric model must be checked. Generic degree and endpoint clause implications, and all recorded cycle cuts, have been checked separately. The reduction from an arbitrary one-singleton vertex to the normalized selection $F$ also remains to be completed. Until these steps are established, the experiment does not prove endpoint proximity or improve the upper bound.

\subsection{An auxiliary conditional singleton estimate}
\begin{proposition}[Conditional estimate from a singleton surplus]
With the quantities in~\eqref{eq:anchored-square-count}, suppose a snake with $v\ge2$ vertices satisfies $T+92\ge2v$. Then it has at most 250 edges.
\end{proposition}
\begin{proof}
The checked incidence identity, $E\ge0$, and $W\le v-2$ give
\[
140v+4T\le36\,856.
\]
Substituting $T\ge2v-92$ yields $148v\le37\,224$. Since $148\cdot252=37\,296$, it follows that $v\le251$ and hence $v-1\le250$. The arithmetic implication has been independently machine-checked; its new aggregate premise remains unproved.
\end{proof}
A Hamming ball of radius two in $\Q_9$ contains $1+9+\binom92=46$ vertices. Thus, if every vertex with exactly one singleton square were proved to lie within distance two of an endpoint, at most 92 vertices could be exceptional. Lemma~\ref{lem:one-square} would then suggest the required aggregate inequality: nonexceptional vertices contribute at least two singleton incidences and exceptional vertices at least one. The endpoint statement and its checked aggregation remain open. This conditional estimate is weaker than the unconditional bounds 245, 241, 238 and 237 proved in this paper. It records the quantitative consequence originally motivating the endpoint-proximity investigation.

\section{Formal verification and further questions}
\subsection{Verification of the mathematical statements}
The explicit 191-edge construction, prefix restriction, singleton-square incidence identity, endpoint and adjacent-count restrictions, exclusion of three consecutive twos, the bound 245, the window exclusions of Section~\ref{sec:window-aggregate}, and the upper bounds 242, 241, 238 and 237 have been checked in Lean 4.19.0. All use the induced-path definition stated in Section~1. The final result has no unresolved proof obligations, additional optimality axioms, or unchecked computational premises. Its logical dependencies are propositional extensionality, classical choice, and quotient soundness, the standard foundations used by the formal library. The complete development behind the bound 241 was built from source, and the dependencies of its final theorems were audited, in two independent computing environments (macOS and Linux). The developments behind the bounds 238 and 237 were built and audited in one environment (Linux), incrementally on top of the already built library; fresh builds from source have not yet been carried out.

Formal verification and computational search serve different purposes here. The finite geometric case analysis is discharged by checked proof terms. The refutations behind the window exclusions for the bounds 241, 238 and 237 are replayed in the kernel by a checker whose soundness is formally proved, so neither the solver nor the clause generator is trusted. Only the conditional bound 239, which the bound 238 subsumes, still rests on window refutations that have been checked only outside Lean. The restricted infeasibility reports in the search section are experimental evidence and have not been converted into independently checked impossibility certificates. The 191-edge witness is a published construction; the elementary finite verification procedure in Appendix~\ref{app:finite-check} also allows its defining conditions to be checked directly from the sequence printed here.

\subsection{A conditional strengthening}
The next proposition isolates a sufficient additional local restriction. It was the first conditional strengthening of the bound 245. Its premise is one of the window exclusions later proved in Section~\ref{sec:window-aggregate}, so its conclusion is superseded by Theorem~\ref{thm:upper241}.

\begin{proposition}[Conditional upper bound of 244]\label{prop:conditional244}
If no snake in $\Q_9$ contains three consecutive internal vertices with singleton-square counts $(2,2,3)$, then every snake in $\Q_9$ has at most 244 edges.
\end{proposition}
\begin{proof}
For a snake with $v\ge2$ vertices, write its singleton counts as $\alpha_0,\ldots,\alpha_{v-1}$ and their sum as $T$. Define
\[
R_i=\begin{cases}
3,&\alpha_i=1,\\
2,&\alpha_i=2,\ i>0,\ \alpha_{i-1}=2,\\
1,&\alpha_i=2\text{ in every other case},\\
0,&\alpha_i\ge3.
\end{cases}
\]
We claim that $5+R_i\le2\alpha_{i+1}+R_{i+1}$ for every $i<v-1$. If $\alpha_{i+1}=1$, its predecessor has count at least four, and both sides equal five. If $\alpha_{i+1}=2$, either its predecessor is at least three and both sides equal five, or its predecessor is two. In the latter case the exclusion of $(2,2,2)$ gives $R_i=1$, $R_{i+1}=2$, so both sides equal six. A predecessor of one is impossible in this case. If $\alpha_{i+1}=3$, a predecessor of one is again impossible; a predecessor of two has $R_i=1$ by the assumed exclusion of $(2,2,3)$, so the left side is at most six. Finally, if $\alpha_{i+1}\ge4$, then $5+R_i\le8\le2\alpha_{i+1}$.

Both endpoint counts are at least four, so $R_0=R_{v-1}=0$ and $2\alpha_0\ge8$. Summing gives
\[
2T\ge8+5(v-1)=5v+3.
\]
The incidence estimate gives $T+35v\le9214$. Therefore $75v+3\le18428$, or $75v\le18425<75\cdot246$. Thus $v\le245$ and the edge count is at most 244. The zero-edge case is immediate.
\end{proof}

The sequence inequality and its conditional numerical consequence have also been checked in Lean. The exclusion of $(2,2,3)$ is one of the window exclusions established formally in Theorem~\ref{thm:upper241}; even the resulting bounds leave a gap between the upper bound and the known construction.

\subsection{Local experiments on the pattern $(2,2,3)$}
A five-vertex induced window has, up to translation and permutation of coordinates, transition sequence $(0,1,2,0)$ or $(0,1,2,3)$. Around its three internal vertices, consider the union of Hamming balls of radius $r$. Require the five window vertices to be selected, selected degree at most two, and degree at least one wherever all nine neighbors are represented. Permit at most two degree-one vertices in this interior, and exclude selected cycles. At each center, its two prescribed neighbors leave 21 eligible opposite square corners; exactly $\alpha$ of these must be absent when its singleton count is $\alpha$.

For prescribed counts $(2,2,3)$, both window shapes admit assignments at radius two. At radius three, the repeated-coordinate model is infeasible using square-cycle exclusions alone; the distinct-coordinate model becomes infeasible after seven additional cycle cuts. Independent checks confirm feasible radius-two assignments and acceptance of restrictions of known snakes. Both radius-three Boolean encodings have contradiction certificates replayed by a separate Python checker. For the repeated-coordinate case, the explicit Boolean contradiction has now been checked in Lean, using a sequence of smaller proofs after the single-file translation exceeded available memory. The distinct-coordinate certificate remains under verification. These results alone do not prove the geometric premise of Proposition~\ref{prop:conditional244}; that premise was later established by a different route, described in Section~\ref{sec:window-aggregate}.

\begin{lemma}[Missing corners at a prescribed internal vertex]\label{lem:window-missing-corners}
Let $S$ be the vertex set of a snake, and suppose $v,v+e_a,v+e_b\in S$, with $a\ne b$. Among the 21 unordered coordinate pairs $\{c,d\}$ disjoint from $\{a,b\}$, exactly $\alpha_S(v)$ have $v+e_c+e_d\notin S$.
\end{lemma}
\begin{proof}
The two specified neighbors exhaust the selected neighbors of $v$, since a vertex of an induced path has degree at most two. A coordinate square using $a$ or $b$ therefore contains another selected vertex and is not a singleton square at $v$. For a pair avoiding $a$ and $b$, both adjacent corners $v+e_c$ and $v+e_d$ are absent. The square is a singleton precisely when its opposite corner is absent. There are $\binom{7}{2}=21$ such pairs, and every singleton square is counted once.
\end{proof}

\begin{proposition}[Exact threshold counters]\label{prop:window-counter}
Let $y_1,\ldots,y_N\in\{0,1\}$ indicate absent corners, and let $k\ge0$ be an integer. Introduce Boolean variables $t_{i,j}$ for $0\le i\le N$ and $0\le j\le k+1$. Set $t_{i,0}$ true and $t_{0,j}$ false for $j>0$. For every $0\le i<N$ and $1\le j\le k+1$, impose the four clauses
\[
\begin{aligned}
&\neg t_{i,j}\lor t_{i+1,j},\\
&\neg y_{i+1}\lor\neg t_{i,j-1}\lor t_{i+1,j},\\
&\neg t_{i+1,j}\lor t_{i,j}\lor y_{i+1},\\
&\neg t_{i+1,j}\lor t_{i,j}\lor t_{i,j-1}.
\end{aligned}
\]
Here $y_i$ is read as true when its value is one. Together with $t_{N,k}$ true and $t_{N,k+1}$ false, these clauses admit a satisfying assignment exactly when $\sum_{i=1}^N y_i=k$.
\end{proposition}
\begin{proof}
Write $C_i=\sum_{h=1}^i y_h$, with $C_0=0$. The four transition clauses are equivalent to
\[
t_{i+1,j}\ \Longleftrightarrow\ t_{i,j}\lor(y_{i+1}\land t_{i,j-1}).
\]
For positive $j$, the integer inequality $C_i+y_{i+1}\ge j$ holds precisely when $C_i\ge j$, or when $y_{i+1}=1$ and $C_i\ge j-1$. Thus the assignment $t_{i,j}=[C_i\ge j]$ satisfies every initial and transition clause. Conversely, induction on $i$ shows that those clauses force this assignment at every represented threshold. The final two clauses are therefore equivalent to $k\le C_N<k+1$, which is $C_N=k$.
\end{proof}

The missing-corner identity, the threshold-counter equivalence, and their connection for each normalized center have been checked in Lean. The finite corner lists have 21 distinct entries and agree with the eligible coordinate pairs. Generic endpoint-count and degree clauses, together with all 112 distinct cycle cuts retained by the two Boolean cores, have also been checked. Completing this particular route would require assembling the exact input clauses under one assignment and establishing the reduction from arbitrary five-vertex windows. Both steps were carried out for the radius-three window model of the next subsection, which proves the exclusion of $(2,2,3)$ among others.

\subsection{The aggregate limit of local windows, and the bounds 242, 241, 238 and 237}\label{sec:window-aggregate}
The same local model can be used for every count pattern. It indicates how far the potential method can go. We decided the radius-three model for every triple of prescribed counts in $\{1,\ldots,6\}^3$ and both window shapes. Of the 216 triples, 144 admit an assignment for at least one shape. Two independent solvers with separate encodings returned the same answer in all 432 cases. Over bi-infinite sequences whose consecutive triples are admitted, the minimum mean count is $14/5$. It is attained by repeating $2,3,3,2,4$. The local restrictions proved above allow mean $7/3$. Larger counts were treated as unrestricted, so these values are lower bounds for the mean.

The minimum $14/5$ is already forced by a small family of patterns. Besides the proved restrictions, it suffices to exclude the internal triples
\[
(2,2,3),\ (2,2,4),\ (2,3,2),\ (3,2,3),\ (4,1,c)\ (4\le c\le6),
\]
together with their reversals. Each of these has, for both window shapes and orientations, an explicit Boolean encoding of the radius-three model. For each encoding a resolution-style refutation was produced and verified by an independent proof checker; these refutations are now also replayed in Lean (Theorem~\ref{thm:upper241}). Restrictions of the 131 bank snakes, the 191-edge witness, and the 137-edge counterexample satisfy every clause of the corresponding encodings. These checks include windows containing counts one and three.

\begin{proposition}[Conditional upper bounds 242 and 241]\label{prop:window-aggregate}
Suppose that no snake in $\Q_9$ contains three consecutive internal vertices whose counts form one of the triples above or their reversals. Then every snake in $\Q_9$ has at most 242 edges. If, in addition, no snake contains four consecutive internal vertices with counts $(2,4,2,3)$ or $(3,2,4,2)$, then every snake has at most 241 edges.
\end{proposition}
\begin{proof}
Let $\beta_0,\ldots,\beta_{n-1}$ be the counts of the internal vertices, so $n=v-2$. Clamp each count above six to seven. Define $P$ on the 49 clamped pairs by the longest-path potential for the weight $14-5z$:
$P(1,4)=3$, $P(2,2)=8$, $P(2,3)=6$, $P(2,4)=3$, $P(3,2)=9$, $P(3,3)=5$, $P(4,1)=12$, $P(4,2)=7$, $P(4,3)=2$, $P(c,1)=9$ and $P(c,2)=4$ for $c\ge5$, and $P=0$ otherwise. A check of all $8^3$ clamped triples shows that
\[
14+P(\beta_{i-1},\beta_i)\le5\beta_{i+1}+P(\beta_i,\beta_{i+1})
\]
for every admitted triple, and that $5(\beta_0+\beta_1)+P(\beta_0,\beta_1)\ge28$. Summing gives $5\sum\beta_i\ge14n-12$. With both endpoint counts at least four, $5T\ge14(v-2)-12+40=14v$. Combined with $T+35v\le9214$ this gives $189v\le46070$, so $v\le243$.

For the second statement, use states consisting of three consecutive clamped counts and the weight $3-w$. The resulting potential is bounded by 4 and verified on all $8^4$ clamped quadruples. It gives $\sum\beta_i\ge3n-4$, so $T\ge3v-6-4+8$. Then $38v\le9216$, and $v\le242$.
\end{proof}

The sequence inequalities, their potentials, and the numerical consequences are checked in Lean for the original snake definition. The proof uses the proved endpoint, count-one, and $(2,2,2)$ results, and keeps the window exclusions as explicit hypotheses. These hypotheses are now formal theorems.

\begin{theorem}[Verified upper bounds 242 and 241]\label{thm:upper241}
Every snake in $\Q_9$ has at most 242 edges, and in fact at most 241 edges. Hence
\[
191\le a(9)\le241.
\]
\end{theorem}
\begin{proof}
By Proposition~\ref{prop:window-aggregate} and Theorem~\ref{thm:lower}, it suffices to establish the window exclusions assumed there. Each exclusion is proved formally in three steps. First, an arbitrary window of a snake with the prescribed counts is moved, by translation and coordinate permutation, to the canonical position of the radius-three local model, and the restriction of the snake is shown to satisfy the defining conditions of that model, including the exact missing-corner counts of Lemma~\ref{lem:window-missing-corners}. Second, every clause of the corresponding unsatisfiable Boolean core is derived from those conditions. Third, the resolution refutation of the core is replayed, segment by segment and by kernel evaluation, through a proof checker whose soundness is itself proved in Lean. The refutation is therefore only data for a verified checker. Combining the exclusions with Proposition~\ref{prop:window-aggregate} gives both bounds.
\end{proof}
The formal proof of Theorem~\ref{thm:upper241} was built from source and audited in two independent environments; its final theorems depend only on propositional extensionality, classical choice and quotient soundness.

All four-vertex windows with counts in $\{1,\ldots,5\}$ were decided. Ignoring which window shape occurs, the minimum mean remains exactly $3$, because $(3,3,3,3)$ is admitted by the shape with transition sequence $(0,1,2,3,0)$. Three consecutive windows of that shape would make vertices seven steps apart adjacent. Tracking transitions, using radius four, or retaining the empty-square term $s_0$ are therefore the natural next refinements.

\paragraph{A direction-aware automaton.} Tracking transitions does help. Let a state consist of the last six transition directions, up to relabelling of coordinates (there are twelve induced shapes), and the last three counts, with counts above six identified (2592 states). A transition appends a direction and a count, and is admitted unless the seven-direction block fails to be an induced path or some window of two, three or four consecutive centres ending at the new vertex is known to be infeasible in the radius-three model. Windows not yet examined are admitted, so the minimum cycle mean of this automaton is a valid lower bound on the average internal count. A refinement loop alternates a minimum-mean-cycle computation with queries to the local model for the windows on that cycle. It stabilises at mean $10/3$ after five rounds; six directions of memory are necessary and sufficient for this value. A dual pruning step reduces the exclusions used to 145 windows, each with an independently checked refutation certificate, and an integer potential of range $8$ on the 2592 states.

\begin{proposition}[Conditional upper bound 239]\label{prop:direction239}
Suppose that no snake in $\Q_9$ contains a window of two, three or four consecutive internal vertices whose transition shape and counts (with counts above six identified) form one of the 145 certified patterns. Then every snake in $\Q_9$ has at most 239 edges.
\end{proposition}
\begin{proof}
With $\beta_j$ the internal counts, the potential gives $3\sum_{j=6}^{m-1}\min(\beta_j,6)\ge10(m-6)-8$. Since $\beta_1,\ldots,\beta_5\ge1$ and both endpoint counts are at least four, $3T\ge10m-36$. Combined with $T+35(m+1)\le9214$ this gives $115m\le27573$, so $m\le239$.
\end{proof}

This statement is checked in Lean for the original snake definition: the step inequality is verified by kernel evaluation over all 877 restricted growth strings of length seven and all count combinations; invariance of the canonical shape under relabelling transports it to arbitrary windows; and the parity formulation of the induced-path condition (a block of flipped coordinates returns to its start exactly when every coordinate appears an even number of times, and ends adjacent to its start exactly when one coordinate appears an odd number of times) connects it to the snake. The 145 window exclusions remain hypotheses. Since the five-centre version below subsumes this bound (Theorem~\ref{thm:upper238}), their formal replay was not completed, and Proposition~\ref{prop:direction239} remains conditional.

With windows of up to five centres and seven directions of memory, the same loop reaches mean $7/2$ with 325 certified exclusions. A further pruning step keeps 175 of them with the same mean, using an integer potential of range $10$ on 44\,064 states.

\begin{proposition}[Conditional upper bound 238]\label{prop:direction238}
Suppose that no snake in $\Q_9$ contains a window of two to five consecutive internal vertices whose transition shape and counts (with counts above six identified) form one of the 175 certified patterns. Then every snake in $\Q_9$ has at most 238 edges.
\end{proposition}
\begin{proof}
The potential gives $7m\le2\sum_{j=1}^{m-1}\alpha_S(v_j)+52$. Adding the endpoint counts, each at least four, gives $2T\ge7m-36$. Combined with $T+35(m+1)\le9214$ this gives $77m\le18394<77\cdot239$, so $m\le238$.
\end{proof}
The corresponding step table is checked by the kernel in 34 pieces, one per canonical seven-direction shape, together with a kernel-checked lemma that every induced window of eight directions extends one of those shapes. The deduction of Proposition~\ref{prop:direction238} from the 175 exclusions has been checked in Lean using only the standard axioms. Each of the 175 exclusions has a refutation certificate verified by an independent checker. These hypotheses are now formal theorems.

\begin{theorem}[Verified upper bound 238]\label{thm:upper238}
Every snake in $\Q_9$ has at most 238 edges. Hence
\[
191\le a(9)\le238.
\]
\end{theorem}
\begin{proof}
By Proposition~\ref{prop:direction238} and Theorem~\ref{thm:lower}, it suffices to establish the 175 window exclusions assumed there. Each is proved formally by the same three steps as in Theorem~\ref{thm:upper241}. An arbitrary window of a snake with the prescribed transition shape and counts is moved by translation and coordinate permutation to the canonical position of the radius-three local model, whose defining conditions the restriction of the snake satisfies; first-occurrence relabelling of the transition directions is realised by a coordinate permutation, and reversed windows are handled through the reversed snake. Every clause of the corresponding unsatisfiable Boolean core is derived from those conditions. Finally a resolution refutation of the core is replayed, segment by segment and by kernel evaluation, through a proof checker whose soundness is proved in Lean. Combining the exclusions with Proposition~\ref{prop:direction238} gives the bound.
\end{proof}
The 175 refutations comprise about $4.6$ million derived clauses with about 270 million resolution hints. Their replay was split into 2\,803 kernel checks, which took about 40 kernel-hours in total, each using at most 4~GB of memory. Because the checker re-verifies every hint, any refutation of the core is acceptable; for 56 of the windows a new refutation was generated from the core itself rather than taken from the original certificate. The final theorem depends only on propositional extensionality, classical choice and quotient soundness. Unlike Theorem~\ref{thm:upper241}, it has so far been audited in a single environment (Linux), against a library built incrementally rather than from source in a fresh directory; an independent rebuild remains to be done.

\paragraph{Windows of radius four.} The local model does not depend on the radius of the ball: a larger radius only adds variables to the Boolean encoding, and the same three steps prove each exclusion. With radius four, five centres and seven directions of memory, the refinement loop reaches mean $42/11$. A cheaper subset of the resulting exclusions, selected by estimated refutation cost, already improves the bound by one: 387 windows give mean $51/14$ with an integer potential of range $61$ on the same states.

\begin{proposition}[Conditional upper bound 237]\label{prop:direction237}
Suppose that no snake in $\Q_9$ contains a window of two to five consecutive internal vertices whose transition shape and counts (with counts above six identified) form one of the 387 certified radius-four patterns. Then every snake in $\Q_9$ has at most 237 edges.
\end{proposition}
\begin{proof}
The potential gives $51m\le14\sum_{j=1}^{m-1}\alpha_S(v_j)+336$. Adding the endpoint counts, each at least four, gives $14T\ge51m-224$. Combined with $T+35(m+1)\le9214$ this gives $541m\le128\,730<541\cdot238$, so $m\le237$.
\end{proof}
The step table is again checked by the kernel in 34 pieces, and the deduction of Proposition~\ref{prop:direction237} from the 387 exclusions has been checked in Lean using only the standard axioms. Each of the 387 exclusions has a refutation certificate verified by an independent checker. These hypotheses are now formal theorems.

\begin{theorem}[Verified upper bound 237]\label{thm:upper237}
Every snake in $\Q_9$ has at most 237 edges. Hence
\[
191\le a(9)\le237.
\]
\end{theorem}
\begin{proof}
By Proposition~\ref{prop:direction237} and Theorem~\ref{thm:lower}, it suffices to establish the 387 window exclusions assumed there. Each is proved formally by the three steps of Theorem~\ref{thm:upper238}: an arbitrary window of a snake with the prescribed transition shape and counts is moved to the canonical position of the radius-four local model, every clause of the corresponding unsatisfiable Boolean core is derived from the conditions of that model, and a resolution refutation of the core is replayed, segment by segment and by kernel evaluation, through the proof checker whose soundness is proved in Lean. Combining the exclusions with Proposition~\ref{prop:direction237} gives the bound.
\end{proof}
The 387 refutations were all generated from the cores themselves; they comprise about $6.4$ million derived clauses with about 289 million resolution hints. Their replay was split into 3\,287 kernel checks, which took about 34 kernel-hours in total. The final theorem depends only on propositional extensionality, classical choice and quotient soundness. Like Theorem~\ref{thm:upper238}, it has so far been audited in a single environment (Linux), against a library built incrementally; an independent rebuild remains to be done.

\paragraph{Limits of local window arguments.} Explicit configurations in $\Q_9$ that are not single paths, such as disjoint unions of induced 26-cycles, satisfy every window condition of the radius-three model with up to five centres. These configurations cap count-and-direction window methods of the type behind Proposition~\ref{prop:direction238} near 227 edges (repository directory \path{milestones/11-local-ceiling/}). Reaching the conjectured value 196 therefore requires global arguments, not only further local window exclusions.

\subsection{A limitation of uniform singleton-density estimates}
The published bank includes a 192-vertex snake for which independent counting gives $T=1632$. This particular census has been checked in Python; its complete connection to the Lean incidence development remains to be checked. It implies that any valid uniform estimate $T\ge cv$ has $c\le17/2$.

The elementary square-census equations alone cannot turn such an estimate into an upper bound below 210 edges. Indeed, the numerical census
\[
(v,a_0,a_1,a_2,a_3)=(211,17,1795,2587,209)
\]
satisfies
\[
a_0+a_1+a_2+a_3=4608,\qquad
 a_1+2a_2+3a_3=36v,\qquad a_3=v-2,
\]
and $a_1\ge(17/2)v$. No path realizing this census is asserted. It demonstrates that these numerical premises permit 211 vertices. The arithmetic implication has been checked in Lean. A stronger conclusion requires additional geometric restrictions, a length-dependent estimate, or another invariant.

\subsection{Overlapping endpoint demands}
Local geometric arguments often produce sets that must contain an endpoint. Their overlap can be retained in the counting argument.
\begin{proposition}[Weighted endpoint accounting]
Let $P$ be a nontrivial snake, and let $(G_i)_{i\in I}$ be a finite family of vertex sets, each containing at least one endpoint of $P$. Assign nonnegative rational weights $w_i$. If
\[
\sum_{i:\,x\in G_i} w_i\le1
\quad\text{for every vertex }x,
\]
then $\sum_{i\in I}w_i\le2$.
\end{proposition}
\begin{proof}
Let $B$ be the two-element endpoint set. Each $G_i$ meets $B$, so
\[
\sum_i w_i\le\sum_i\sum_{x\in B\cap G_i}w_i
=\sum_{x\in B}\sum_{i:\,x\in G_i}w_i\le2.
\]
The integer-weight version, obtained by clearing a common denominator, has been checked in Lean.
\end{proof}

For example, suppose five distinct candidate vertices $u_0,\ldots,u_4$ yield demands $G_i=\{u_i,u_{i+1}\}$, with subscripts taken modulo five. Giving each demand weight $1/2$ satisfies the per-vertex condition, but the total weight is $5/2$, a contradiction. No three of these pairs are disjoint, so a test limited to three disjoint endpoint demands would miss this obstruction. The five-pair contradiction is also checked independently in Lean. Identifying a forced occurrence of this arrangement, or another weighted obstruction, in low-singleton geometry is an open research direction. The counting principle is standard; no claim of historical novelty is made for it.

\subsection{Conditional share targets and a finite relaxation obstruction}
Two empty-square charges must be distinguished. For an empty-square corner $q$,
let $e(q)$ count empty coordinate squares containing it and let $d(q)$ count its
selected neighbours. The integer corner share received by a selected vertex $x$
at resolution $M$ is
\[
s_M(x)=\sum_{q\sim x}\left\lfloor\frac{M e(q)}{d(q)}\right\rfloor,
\]
with division by zero assigned the value zero. Double counting and the division
algorithm give the exact identity
\[
\sum_{x\in S}s_M(x)+\sum_q (M e(q)\bmod d(q))=M E,
\]
where $E=4s_0$ and $n\bmod0=n$. Thus undominated empty corners and integer rounding
are both retained in the loss term. For $M=12$ put
$G=\sum_{x\in S}(24\alpha_S(x)+s_{12}(x))$. For a nontrivial snake the incidence
bound implies $G+840v\le221136$. Consequently $G\ge114v-1145$ would imply at most
231 edges, and $G\ge114v-191$ would imply at most 230 edges: the excluded vertex
counts would require respectively $954\cdot233\le222281$ and
$954\cdot232\le221327$, both false by exactly one integer unit.
These lower bounds for $G$ have not been proved for arbitrary snakes.

Checking a few favourable periodic patterns is insufficient to prove such a
lower bound. A finite abstract relaxation with 232 edges and 233 vertices has
endpoint counts four, internal counts four except for the 46 indices divisible
by five and index 116, which have count five. It has $T=979$, $E=160$, and may
assign every corner share zero, giving $G=23496<25417$. It satisfies the existing
five-centre exclusion and saved packing constraints. Its transition word repeats
$0,1,\ldots,8$, so vertex 18 repeats vertex zero: it is not a snake. It identifies
a missing restriction in this abstraction rather than a counterexample to the
desired geometric theorem.

A different charge assigns weights $6,3,2$ across parallel squares having one,
two, or three selected corners. If $A$ is the total internal singleton count and
$H$ is this total parallel-square share, its global bound is
$21(A+8)+H+735(m+1)\le193494$. The proposed six-centre premises
$9m\le2A+59$ and $21(m-7)\le6H$ then imply
$1666m\le386470$, hence $m\le231$. For seven centres, premises
$89m\le19A+C$ and $27(m-8)\le7H$ imply
$111351m\le25618707+147C$; $C\le703$ suffices for 230 edges.
The arithmetic and elementary share-gate identities have been checked in Lean.
The complete Boolean refutations and their applicability to all snake windows
remain required. The two share definitions and their certificates are not interchangeable.

The earlier 237 bound retains its contributor audit. A fresh independent rebuild
of its complete custom source dependencies has not been completed in this
investigation; generated refutation modules omitted from the source archive
must first be reconstructed. Neither 231 nor 230 is asserted as an unconditional
theorem here.

\subsection{The exact optimum}
The unconditional conclusion is $191\le a(9)\le237$ (Theorem~\ref{thm:upper237}), which subsumes the earlier verified bounds 242, 241 and 238 of Theorems~\ref{thm:upper241} and~\ref{thm:upper238} and the conditional bound 239 of Proposition~\ref{prop:direction239}. The full radius-four window set of the refinement loop, 732 windows with mean $42/11$, would give 236 edges by the same arithmetic. Its refutation certificates have been checked outside Lean, the bound has been compiled in Lean conditionally on the window exclusions, and the kernel replays of the 361 refutations not shared with the 237 set are in progress. The bound 236 is not established yet. By the limits described in Section~\ref{sec:window-aggregate}, local window arguments of this kind cannot reach much below 227 edges.
No 192-edge construction or universal exclusion of length 192 has been obtained. A solution requires a length $L$ for which both an explicit snake and an exclusion of length $L+1$ are established; the prefix argument then excludes all longer paths. Conjecture~\ref{conj:196} would require a 196-edge construction and an exclusion of length 197, but the proved bounds do not single out that value. Further progress may come from stronger restrictions on singleton counts or from constructions outside the tested families.

\appendix
\section{Explicit transition sequences and blockers}\label{app:witness}
All coordinates are numbered from zero. Spaces only separate groups; start at the zero vertex unless another initial vertex is stated, and apply each coordinate flip in order.
\subsection{Verified 191-edge nine-dimensional witness}
\begin{quote}\small\ttfamily
0120314506 2106502631 4321035023 6210623607 4063260126\\
3206483026 3153263120 3623402362 0162360750 6326012632\\
0652362013 4120132610 2156210263 0210674601 2063201260\\
4816201351 2013261021 4621026302 1067560120 6
\end{quote}
Its last vertex is 118. For moves in coordinates $0,\ldots,8$, respectively, earlier blocking indices are
\[
(19,37,13,13,27,19,189,43,91).
\]
For each coordinate $k$, put $w=v_{191}\mathbin{\oplus}2^k$. Direct calculation gives $w=v_i$ or $\Ham(w,v_i)=1$ for the corresponding displayed index $i<191$.

\subsection{Verified 93-edge endpoint-maximal counterexample}
\begin{quote}\small\ttfamily
8624301806 7812851450 7860736287 5286237138 0746235243\\
8145783485 6180720370 4716078615 6326120174 036
\end{quote}
The terminal vertex is 289. Blocking indices for its nine moves are
\[
(1,33,51,91,45,1,91,81,45).
\]
For the reversed ordering, whose terminal vertex is zero, the blocking indices in that reversed ordering are
\[
(25,19,43,27,19,61,21,23,91).
\]
The reversed ordering is $v_{93},v_{92},\ldots,v_0$. Its pairwise distances are unchanged. Each listed blocker precedes the reversed endpoint and either equals the candidate next vertex or has Hamming distance one from it.

\subsection{Verified 98-edge eight-dimensional seed}
\begin{quote}\small\ttfamily
0123104154 3146340134 6743104274 0147514016 3104154314\\
2340134274 3104674014 7561340142 7410431641 04764013
\end{quote}
The seed is attributed to Carlson and Hougen by the source page \cite{records}. Each transition sequence in this appendix determines every vertex by the recurrence stated at the start of the paper.

\subsection{A 137-edge counterexample for singleton squares}\label{app:square-counterexample}
Start at the integer vertex 145, rather than zero, and use the following coordinate transitions:
\begin{quote}\small\ttfamily
0132518601 5823563751 0231051201 8510231061 2714831024\\
7201352471 3814601234 6187510231 0512018510 2376320158\\
1021501320 1871023158 1321531023 1581320
\end{quote}
The resulting ordered path has 138 distinct vertices. Its vertex at index 127 is zero, and its last vertex is 288. At zero, the only singleton coordinate square uses coordinates 2 and 3.

\subsection{Compatible paths of lengths 96 and 94}\label{app:joint-seed}
Both projected paths start at zero in $\Q_8$. For the first path use
\begin{quote}\small\ttfamily
0213047013 0412031052 1031243103 2672301324 1301237103\\
1240541240 1342147314 2130214236 0135231421 032142
\end{quote}
For the second path use
\begin{quote}\small\ttfamily
1031240741 2401342145 3142130214 2367324123 0124132714\\
2130450130 4120310721 0312431032 6523013241 3012
\end{quote}
The projected vertex sets have intersection $\{0\}$ only. Reverse the first sequence of vertices, place it in layer zero, cross at zero, and follow the second path in layer one.


\section{Finite verification of the printed witnesses}\label{app:finite-check}
For a transition sequence $t_1,\ldots,t_L$ in dimension $n$ and initial vertex $z$, the following finite procedure verifies the snake conditions without any additional data.
\begin{enumerate}
\item Check $0\le z<2^n$ and $0\le t_k<n$ for every $k$. Starting with $v_0=z$, compute $v_k=v_{k-1}\oplus2^{t_k}$ for $1\le k\le L$.
\item For every $0\le i<j\le L$, compute the number of nonzero binary digits of $v_i\oplus v_j$. Reject if this number is zero. Also reject if it equals one and $j\ne i+1$, or if it differs from one and $j=i+1$.
\item Accept if no pair is rejected. The reconstructed sequence is then an induced path of exactly $L$ edges by definition.
\end{enumerate}
For a displayed endpoint blocker $i$ in coordinate $k$, compute $w=v_L\oplus2^k$ and check $i<L$ and either $w=v_i$ or $\Ham(w,v_i)=1$. Apply the same calculation to the reversed vertex ordering for left-endpoint blocking. For the two compatible layer paths, verify each path separately and check that their projected vertex sets intersect only at zero.

This procedure requires $\binom{L+1}{2}$ pair tests for a single witness. It verifies the explicit constructions and counterexamples; it does not test all possible snakes and therefore supplies no global upper bound.

\subsection{A radius-three linear forest}\label{app:radius-three}
The following lists give vertices as decimal integers encoding nine-bit words. Each row is one path; the rows are distinct connected components of the induced graph on their union.
\begin{align*}
P_{1}&=(134,132,140,136,138,10,2,0,1,5,261,260,268,264,265)\\
P_{2}&=(22,20,28,24,25)\\
P_{3}&=(38,36,44,40,41)\\
P_{4}&=(49,48,50)\\
P_{5}&=(70,68,76,72,73)\\
P_{6}&=(81,80,82)\\
P_{7}&=(97,96,98)\\
P_{8}&=(145,144,146)\\
P_{9}&=(161,160,162)\\
P_{10}&=(193,192,194)\\
P_{11}&=(273,272,274)\\
P_{12}&=(289,288,290)\\
P_{13}&=(321,320,322)\\
P_{14}&=(385,384,386)
\end{align*}
The union has 60 vertices, including zero. It is a linear forest with 14 components, not a single snake.

\clearpage
\addcontentsline{toc}{section}{References}
\begin{thebibliography}{99}\raggedright
\bibitem{oeis} OEIS Foundation. \emph{A099155: Maximum length of a simple path with no chords in the $n$-dimensional hypercube}. \url{https://oeis.org/A099155}. Accessed 13 October 2026.
\bibitem{ostergard} P.~R.~J. Östergård and V.~H. Pettersson. Exhaustive Search for Snake-in-the-Box Codes. \emph{Graphs and Combinatorics} 31 (2015), 1019--1028. \url{https://doi.org/10.1007/s00373-014-1423-3}.
\bibitem{orland} P.~Orland et al. \emph{New Snake-in-the-Box Records via Snakepit Surgery and Learned Construction}. arXiv:2607.15270v3, 30 September 2026. \url{https://arxiv.org/html/2607.15270v3}.
\bibitem{data} Math-AI-Caltech. \emph{Snake-in-the-Box: public record data}. \url{https://github.com/Math-AI-Caltech/Snake-in-the-Box}. Accessed 15 October 2026.
\bibitem{records} T.~Ace. \emph{Snake-in-the-Box lower bounds and transition sequences}. \url{https://www.minortriad.com/snake/}. Accessed 13 October 2026.
\bibitem{taylor} T.~Taylor. \emph{Constructing longer snakes and improved asymptotic bounds in hypercubes}. arXiv:2609.27483v1, 23 September 2026. \url{https://arxiv.org/html/2609.27483v1}.
\bibitem{huang} H.~Huang. Induced subgraphs of hypercubes and a proof of the Sensitivity Conjecture. \emph{Annals of Mathematics} 190 (2019), 949--955. \url{https://annals.math.princeton.edu/2019/190-3/p06}.
\bibitem{lukito} A.~Lukito and A.~J. van Zanten. A New Non-Asymptotic Upper Bound for Snake-in-the-Box Codes. \emph{Journal of Combinatorial Mathematics and Combinatorial Computing} 39 (2001), 147--156. \url{https://combinatorialpress.com/article/jcmcc/Volume%20039/vol-039-paper%2012.pdf}.
\bibitem{danzerklee} L.~Danzer and V.~Klee. Lengths of Snakes in Boxes. \emph{Journal of Combinatorial Theory} 2 (1967), 258--265. \url{https://doi.org/10.1016/S0021-9800(67)80026-7}.
\bibitem{douglas} R.~J. Douglas. Upper Bounds on the Length of Circuits of Even Spread in the $d$-Cube. \emph{Journal of Combinatorial Theory} 7 (1969), 206--214. \url{https://doi.org/10.1016/S0021-9800(69)80013-X}.
\bibitem{deimer} K.~Deimer. A New Upper Bound for the Length of Snakes. \emph{Combinatorica} 5 (1985), 109--120. \url{https://doi.org/10.1007/BF02579373}.
\bibitem{soloveva} F.~I. Solov'eva. An Upper Bound on the Length of a Cycle in the $n$-Dimensional Unit Cube (in Russian). \emph{Metody Diskretnogo Analiza} 45 (1987), 71--76.
\bibitem{emelyanov} P.~G. Emel'yanov. On an Upper Bound for the Length of a Snake in an $n$-Dimensional Unit Cube (in Russian). \emph{Diskretnyi Analiz i Issledovanie Operatsii} 2(3) (1995), 10--17. English translation in \emph{Operations Research and Discrete Analysis}, Kluwer, 1997. \url{https://doi.org/10.1007/978-94-011-5678-3_4}.
\bibitem{snevily} H.~S. Snevily. The Snake-in-the-Box Problem: A New Upper Bound. \emph{Discrete Mathematics} 133 (1994), 307--314.
\bibitem{zemor} G.~Zémor. An Upper Bound on the Size of the Snake-in-the-Box. \emph{Combinatorica} 17 (1997), 287--298. \url{https://doi.org/10.1007/BF01200911}.
\bibitem{emelyanovlukito} P.~G. Emelyanov and A.~Lukito. On the Maximal Length of a Snake in Hypercubes of Small Dimension. \emph{Discrete Mathematics} 218 (2000), 51--59.
\bibitem{wynn} E.~Wynn. \emph{Constructing Circuit Codes by Permuting Initial Sequences}. arXiv:1201.1647, 2012. \url{https://arxiv.org/abs/1201.1647}.
\end{thebibliography}
\end{document}
